%=====================================================================
% 光学（02）外教 Lecture 4 · Electro-optic & Magneto-optic Effects
% 精讲讲义（中文讲解 + LaTeX 原生公式）
% 来源：subjects/02-光学/课件/Electro-optic & Magneto-optic Effects.pdf（82 页）
%       作业：subjects/02-光学/课件/Homework EO MO.docx（= 课件 p82）
% 数值核对：subjects/02-光学/讲义/eomo-check/verify_l4.py（42 条断言）
% 编译：xelatex -halt-on-error 两遍
%=====================================================================
\documentclass[11pt]{ctexart}
%=====================================================================
% theme-violet.tex —— 讲义精讲 · 紫花主题
% 用法：\documentclass[11pt]{ctexart}
%       \input{assets/theme-violet}
%       \begin{document} ... \end{document}
% 只用 ctexart 自带能力 + xcolor/enumitem/fancyhdr/pifont/amsmath，
% 不依赖 tikz/tcolorbox，保证 MiKTeX 可直接编译。
%=====================================================================

\usepackage[a4paper,left=2.0cm,right=2.0cm,top=2.1cm,bottom=2.3cm,%
            headsep=12pt,footskip=22pt]{geometry}
\usepackage{amsmath,amssymb}
\usepackage{needspace}
\usepackage{xcolor}
\usepackage{array}
\usepackage{enumitem}
\usepackage{pifont}
\usepackage{fancyhdr}
\usepackage[hidelinks]{hyperref}

%---------------------------------------------------------------- 配色
\definecolor{violet}{RGB}{107,78,143}      % 主色：紫
\definecolor{violetdeep}{RGB}{63,44,92}    % 深紫：标题/强调
\definecolor{petal}{RGB}{183,159,208}      % 花瓣：装饰线、花饰
\definecolor{bandviolet}{RGB}{247,243,251} % 浅紫底：档案/提示框
\definecolor{bandcream}{RGB}{252,249,242}  % 米白底：易错/注意框
\definecolor{rulepurple}{RGB}{217,207,230} % 分隔线
\definecolor{inkgray}{RGB}{92,96,110}      % 次要文字

%---------------------------------------------------------------- 字体
% 有 Noto CJK 用 Noto，没有就退回 ctex 默认（Windows 下为宋体/黑体）
\IfFontExistsTF{Noto Serif CJK SC}{\setCJKmainfont{Noto Serif CJK SC}}{}
\IfFontExistsTF{Noto Sans CJK SC}{\setCJKsansfont{Noto Sans CJK SC}}{}

%---------------------------------------------------------------- 版式
\setlist{nosep,leftmargin=1.9em}
\linespread{1.28}
\setlength{\parskip}{3.5pt}
\setlength{\parindent}{0pt}
\setlength{\tabcolsep}{5pt}
\renewcommand{\arraystretch}{1.28}

%---------------------------------------------------------------- 花饰
% \ding{95} = 五瓣小花，\ding{96} = 星形花；都来自 pifont，必定存在
\newcommand{\orn}{\textcolor{petal}{\ding{95}}}          % 行内小标
\newcommand{\ornstar}{\textcolor{petal}{\ding{96}}}      % 星形花
\newcommand{\flowerrule}{%                                % 居中花饰分隔线
  \par\vspace{7pt}%
  \noindent\makebox[\linewidth]{\color{petal}%
    \rule{0.20\linewidth}{0.6pt}\hspace{8pt}{\small\ding{95}}\hspace{4pt}%
    \raisebox{1.1pt}{\tiny\ding{95}}\hspace{4pt}{\small\ding{95}}\hspace{8pt}%
    \rule{0.20\linewidth}{0.6pt}}%
  \par\vspace{7pt}}

%---------------------------------------------------------------- 标题与节头
\newcommand{\doctitle}[3]{%   #1 主标题  #2 副标题  #3 元信息（可用 \\ 分行）
  \begin{center}
    {\color{petal}\small \orn\enspace #2\enspace \ornstar}\\[7pt]
    {\LARGE\bfseries\color{violetdeep} #1}\\[10pt]
    {\color{petal}\rule{0.55\linewidth}{1pt}}\\[4pt]
    {\small\color{inkgray} #3}
  \end{center}\vspace{2pt}}

\newcommand{\hdr}[2]{%   #1 章节号（如 §0）  #2 标题
  \par\vspace{12pt}%
  \noindent\colorbox{violetdeep}{\color{white}\bfseries\small\;#1\;}\hspace{7pt}%
  {\color{violetdeep}\large\bfseries #2}\par\nobreak
  \vspace{3pt}{\color{petal}\hrule height 1.2pt}\vspace{6pt}}

\newcommand{\subhdr}[1]{%
  \par\vspace{6pt}%
  \noindent{\color{violet}\bfseries #1}\par\vspace{1pt}}

%---------------------------------------------------------------- 提示框
\newcommand{\callout}[1]{% 浅紫底：讲义档案、读法说明
  \par\medskip\noindent
  \fcolorbox{rulepurple}{bandviolet}{\parbox{\dimexpr\linewidth-2\fboxsep-2\fboxrule}{\vspace{2pt}#1\vspace{2pt}}}%
  \par\medskip}

\newcommand{\warnbox}[1]{% 米白底：易错、边界条件
  \par\medskip\noindent
  \fcolorbox{rulepurple}{bandcream}{\parbox{\dimexpr\linewidth-2\fboxsep-2\fboxrule}{\vspace{2pt}#1\vspace{2pt}}}%
  \par\medskip}

\newcommand{\gapitem}[2]{% #1 缺口类型（紫） #2 内容
  \par\vspace{2pt}\noindent
  \hangindent=1.4em\hangafter=1
  {\color{petal}\small\ding{96}}\enspace{\bfseries\color{violetdeep}#1}\enspace #2\par}

% 题干块（\qstem）：每道例题／作业的「详解」开头必须先复述题干与出处。
% 起因：2026-10-25 第五章讲义把七道作业的题面集中放在一页、详解另起一页，
% 读者翻到「第 4 题」时看不到条件，等于一道没有题干的题目。
% 规则：详解块内自带题干，读者不翻页就能读懂题在问什么。
\newcommand{\qstem}[1]{%
  \par\medskip\noindent
  \fcolorbox{violet}{bandviolet}{%
    \parbox{\dimexpr\linewidth-2\fboxsep-2\fboxrule}{%
      \vspace{2pt}%
      {\small\bfseries\color{violet}题干}\hspace{0.8em}#1%
      \vspace{2pt}}}%
  \par\medskip}

% v2: 保护题干与开头，长题可跨页，续页补上下文。
\newcommand{\qbreak}{\needspace{6\baselineskip}}
% 独立练习：ID、来源、题干、考点、答案标签（通常 ans:ID）。
\newcommand{\practice}[5]{%
  \qbreak\subhdr{练习 #1}\qstem{（#2）#3}%
  \textbf{考点}：#4\par
  详解见附录第~\pageref{#5}~页。\par\medskip}
% 附录开头：ID、来源、条件摘要；自动产生 ans:ID 标签。
\newcommand{\answerstart}[3]{%
  \qbreak\subhdr{#1\quad 详解}\label{ans:#1}\qstem{（#2）#3}}
\newcommand{\answercontinue}[2]{%
  \needspace{5\baselineskip}\subhdr{#1\quad 续解}%
  {\small\color{inkgray}本页接续：#2}\par\medskip}
% 延迟提供默认值，让旧讲义可在导言区保留自己的 \newcommand{\step}。
\AtBeginDocument{\providecommand{\step}[1]{\needspace{3\baselineskip}\par\medskip\textbf{#1}\par}}

% 前提旁注：只放在「需要理解的那一步」后面，给陈述 + 本题代入，不重证。
\newcommand{\aside}[1]{%
  \par\vspace{1pt}\noindent
  \hangindent=1.6em\hangafter=1
  {\small\color{inkgray}%
    \hspace*{0.6em}{\color{petal}$\triangleright$}\hspace{2pt}#1}%
  \par\vspace{2pt}}

%---------------------------------------------------------------- 表格
% 本机 MiKTeX 的 colortbl 与 array 版本不匹配（\rowcolor 会报错），
% 故表格统一用「上粗紫线 + 表头 + 中细线」的三线式，不依赖 colortbl。
\newcommand{\tabtop}{\par\vspace{3pt}{\color{violet}\hrule height 1.2pt}\vspace{1pt}}
\newcommand{\tabmid}{\par\vspace{1pt}{\color{petal}\hrule height 0.5pt}\vspace{1pt}}
\newcommand{\tabbot}{\par\vspace{3pt}{\color{rulepurple}\hrule height 0.5pt}\vspace{8pt}}
\newcommand{\thead}[1]{{\bfseries\color{violetdeep}#1}}

%---------------------------------------------------------------- 页眉页脚
\pagestyle{fancy}
\fancyhf{}
\renewcommand{\headrulewidth}{0.4pt}
\renewcommand{\headrule}{\hbox to\headwidth{\color{rulepurple}\leaders\hrule height \headrulewidth\hfill}}
\fancyhead[L]{\small\color{inkgray}\songti \doctag}
\fancyhead[R]{\small\color{inkgray}\songti \coursetag}
\fancyfoot[C]{\small\color{inkgray}\orn\enspace{\color{violet}\thepage}\enspace\orn}

% 每篇讲义在导言区设这两个值
\newcommand{\doctag}{}
\newcommand{\coursetag}{}
\newcommand{\settags}[2]{\renewcommand{\doctag}{#1}\renewcommand{\coursetag}{#2}}

%---------------------------------------------------------------- 数学记号
\newcommand{\vect}[1]{\boldsymbol{#1}}
\newcommand{\dd}{\,\mathrm{d}}
\newcommand{\ddx}{\frac{\mathrm{d}}{\mathrm{d}x}}
\addtolength{\jot}{2pt}



\usepackage{tikz}
\usetikzlibrary{arrows.meta,calc,positioning,decorations.pathmorphing}
\usepackage{graphicx}
\setlength{\headheight}{22pt}
% theme 的 \tabmid 在 tabular 内会落进首格，改成 \noalign 形式
\renewcommand{\tabmid}{\noalign{\vskip 2pt{\color{petal}\hrule height 0.5pt}\vskip 3pt}}
\settags{Electro-optic \& Magneto-optic · L4 精讲}{光学（02）· 外教 Lecture 4}
\let\subhdrtheme\subhdr
\renewcommand{\subhdr}[1]{\needspace{3\baselineskip}\subhdrtheme{#1}}
\let\hdrorig\hdr
\renewcommand{\hdr}[2]{\needspace{8\baselineskip}\hdrorig{#1}{#2}}
\setlength{\emergencystretch}{1.6em}
\tikzset{>={Stealth[length=2.0mm]},
  box/.style={draw=violet,rounded corners=2pt,fill=bandviolet,align=center,inner sep=3.5pt},
  boxc/.style={draw=petal,rounded corners=2pt,fill=bandcream,align=center,inner sep=3.5pt},
  bx/.style={draw=violet,thick,align=center,inner sep=2pt}}
% Verdet 常数用花体 V，避免与电压 V 撞符号
\newcommand{\vdt}{\mathcal{V}}
% 主题只给了 \callout{...}、\warnbox{...} 的命令形式，不能容纳分段、列表与行间公式。
% 这里补两个视觉等价的环境（minipage + fcolorbox），供较长的提示框使用。
\newsavebox{\dshbox}
\newenvironment{Callout}{%
  \par\medskip\noindent
  \begin{lrbox}{\dshbox}\begin{minipage}{\dimexpr\linewidth-2\fboxsep-2\fboxrule\relax}%
}{%
  \end{minipage}\end{lrbox}%
  \fcolorbox{rulepurple}{bandviolet}{\usebox{\dshbox}}\par\medskip}
\newenvironment{WarnBox}{%
  \par\medskip\noindent
  \begin{lrbox}{\dshbox}\begin{minipage}{\dimexpr\linewidth-2\fboxsep-2\fboxrule\relax}%
}{%
  \end{minipage}\end{lrbox}%
  \fcolorbox{rulepurple}{bandcream}{\usebox{\dshbox}}\par\medskip}

\begin{document}

\doctitle{电光与磁光效应（第四讲）}
{普克尔斯盒 · 克尔效应 · 法拉第旋转 · 光隔离器}
{依据：\texttt{Electro-optic \& Magneto-optic Effects.pdf}（82 页，Lecture 4）\\
作业：课件 p82 三题 —— 题干见 §16，中文详解见附录 C，英文答题版见附录 D\\
光学（02）｜整理日期 2026-11-20}


%=====================================================================
\hdr{§0}{这一讲要解决什么问题}

前三讲分别是：怎么描述偏振（L1）、怎么算偏振态的变化（L2）、怎么用器件造和测偏振（L3）。
这一讲把「外场」加进来：\textbf{不碰器件，只拧一个电压或一个磁场，光的偏振和强度就跟着变}。

整堂课压成三句话：

\textbf{一、场改的是折射率，不是光。} Kerr 和 Pockels 效应都是电场改了晶体的折射率椭球，
于是两个正交偏振分量之间多出一个相位差 $\Gamma$；磁光效应同理，只是它把右旋和左旋圆偏振分开，
改的是偏振面的方位。两者最后都落到同一个算式 $\Gamma=2\pi\Delta n\,L/\lambda$（§1）。

\textbf{二、对称性决定「线性还是二次」。}有反演中心的材料，$\pm E$ 必须给同一个结果，
所以奇次项被禁掉：只剩 $\Delta n\propto E^2$（Kerr）；没有反演中心才有 $\Delta n\propto E$（Pockels）。
这一条判据能让你在没查表时就说出「这块晶体能不能做普克尔斯盒」（§2）。

\textbf{三、同一句话里的两个方向，一个有、一个没有。}把一束光送进去再原路送回来：
Faraday 旋转会继续同向累加（$2\theta$），自然旋光会原路退回去（$0$）。
这个「非互易」是光隔离器能单向通光的全部理由，也是本讲作业第 2 题在问的事（§9、§12）。

\begin{center}
\begin{tikzpicture}[font=\footnotesize,x=1cm,y=1cm]
\node[box,text width=13.2cm] (q) at (7.4,1.0)
  {\textbf{核心问题：不换光学元件，只加外场 $\boldsymbol{E}$ 或 $\boldsymbol{B}$，怎么让光的偏振面、相位、强度跟着变？}};

\node[box,text width=4.2cm] (e0) at (3.0,-0.9) {\textbf{电场支路}\\ 场改折射率椭球};
\node[box,text width=4.2cm] (b0) at (11.4,-0.9) {\textbf{磁场支路}\\ 场把 RCP 与 LCP 分开};

\node[box,text width=4.2cm] (e1) at (3.0,-2.4) {$\Delta n\propto E^2$：\textbf{Kerr}（§3）\\ 所有材料，弱};
\node[box,text width=4.2cm] (e2) at (3.0,-3.7) {$\Delta n\propto E$：\textbf{Pockels}（§5--§6）\\ 仅非中心对称，强};
\node[box,text width=4.2cm] (b1) at (11.4,-2.4) {$\theta=\vdt BL$：\textbf{Faraday}（§11）\\ 非互易：往返 $2\theta$};
\node[box,text width=4.2cm] (b2) at (11.4,-3.7) {反射变：\textbf{MOKE}（§13）\\ 线双折射：Voigt（§14）};

\node[boxc,text width=9.6cm] (m) at (7.4,-5.2)
  {两条路落到同一把尺子：$\Gamma=\dfrac{2\pi}{\lambda}\Delta n\,L$（线性相位），或 $\theta$（方位角）};
\node[boxc,text width=9.6cm] (m2) at (7.4,-6.4)
  {正交检偏器读出强度：$T=\sin^{2}\dfrac{\Gamma}{2}=\sin^{2}\dfrac{\pi V}{2V_{\pi}}$（§7）};

\node[box,text width=4.6cm] (d1) at (3.0,-7.9) {\textbf{Pockels 盒 / Kerr 盒}\\ 强度调制器 · 相位调制器 · Q 开关};
\node[box,text width=4.6cm] (d2) at (11.4,-7.9) {\textbf{Faraday 旋转器}\\ 光隔离器 · 环形器 · 磁光存储};

\node[boxc,text width=13.0cm] (k) at (7.4,-9.35)
  {\textbf{一条判据贯穿全讲}：自然旋光（互易，往返 $0^\circ$）\\ vs.\ Faraday 旋转（非互易，往返 $2\theta$）——§9 与 §12 专门对照};

\draw[->,violet] (q.south) -- ++(0,-0.35) -| (e0.north);
\draw[->,violet] (q.south) -- ++(0,-0.35) -| (b0.north);
\draw[->,violet] (e0)--(e1); \draw[->,violet] (e1)--(e2);
\draw[->,violet] (b0)--(b1); \draw[->,violet] (b1)--(b2);
\draw[->,violet] (e2.south) -- ++(0,-0.4) -| (m.north);
\draw[->,violet] (b2.south) -- ++(0,-0.4) -| (m.north);
\draw[->,violet] (m)--(m2);
\draw[->,violet] (m2.south) -- ++(0,-0.35) -| (d1.north);
\draw[->,violet] (m2.south) -- ++(0,-0.35) -| (d2.north);
\draw[->,violet] (3.0,-8.45) -- (3.0,-8.85);
\draw[->,violet] (11.4,-8.45) -- (11.4,-8.85);
\node[text width=15.4cm,align=left,font=\scriptsize,text=inkgray] at (7.4,-11.0)
 {读法：上排是「场怎么起作用」，中排是两条路共用的换算（先得到相位差或方位角），下排才是器件。\\
  左边一支改的是\emph{线}双折射（两个正交线偏振的相位差），右边一支改的是\emph{圆}双折射（RCP 与 LCP 的相位差）——
  这个区别决定了最下排那条互易性判据。};
\end{tikzpicture}
\end{center}

\subhdr{0.1 页码路线：82 页课件对应到本讲义}

{\small\color{inkgray}课件 p82 就是本次作业题面；p43--p52、p79--p81 是课件自带的三加三道例题，
本讲义 §8 把它们当\emph{教学示范}讲（课件已给答案，不再作为独立练习）。}

\begin{center}
\small
\tabtop
\begin{tabular}{@{}p{1.9cm}>{\raggedright}p{8.7cm}>{\raggedright\arraybackslash}p{4.6cm}@{}}
\thead{课件页} & \thead{内容} & \thead{本讲义} \\
\tabmid
p1--p4 & 标题、计划与学习目标、为什么重要、效应分类表 & §0、§1 \\
\tabmid
p5--p6 & 偏振态与琼斯矢量、折射率与双折射、指标体、$\Gamma=2\pi\Delta nL/\lambda$ & §1.3--§1.5 \\
\tabmid
p7--p8 & 为什么有的线性有的二次（对称性）、材料类型表 & §2 \\
\tabmid
p9--p17 & Kerr：小史、$\Delta n=K\lambda E^2$、两种机理、常数表、温度与波长、液固对比、盒几何、响应时间 & §3 \\
\tabmid
p18--p20 & Kerr 三问 FAQ（反极性、为何还用它、光能否平行场） & §3.7 \\
\tabmid
p21--p32 & Pockels：小史、$\Delta n=-\tfrac12n^3rE$、对称性表、中心对称为何无线性项、张量 $r_{ij}$、纵横向、装置图、$V_\pi$、材料表、强度调制、对比表、工程难点 & §5--§7 \\
\tabmid
p33--p35 & Pockels 三问 FAQ（原子层面、纵横选型、速度极限） & §6.5、§7.7 \\
\tabmid
p36--p42 & 应用目录、通用振幅调制器、调制曲线、三个工作点、线性区、相位调制器与 MZI & §7 \\
\tabmid
p43--p52 & 例题 1（KDP 横向 $V_\pi$）、例题 2（Kerr 盒电压）、例题 3（强度调制） & §8 \\
\tabmid
p53--p59 & 磁光引言、MCB/MCD 与非互易、发现年表、名词表、介电张量、Faraday 公式、物理机制 & §10--§11 \\
\tabmid
p60--p65 & MOKE：原理、三种几何表、极向、纵向、横向 & §13 \\
\tabmid
p66--p69 & Voigt 与 Cotton--Mouton、磁性类型与磁光响应、磁光材料分类 & §14.1、§14.4、§15.4 \\
\tabmid
p70--p73 & 圆二色性、CD 的量、CD 光谱仪、CB 与 CD 对照 & §14.2--§14.3 \\
\tabmid
p74--p78 & 隔离器与环形器、磁光存储、三问 FAQ（非互易、Faraday vs MOKE、为何弱） & §12.4、§15 \\
\tabmid
p79--p81 & 磁光三道例题（TGG 转角、隔离器长度、两材料比较） & §8.4 \\
\tabmid
p82 & \textbf{本次作业三题} & §16、附录 C、附录 D \\
\end{tabular}
\tabbot
\normalsize
\end{center}

\subhdr{0.2 记号（全篇统一）}

{\small\color{inkgray}课件把一个 $V$ 同时用作「电压」和「Verdet 常数」。本讲义把 Verdet 常数写成花体
$\vdt$，电压一律写 $V$；作业里若照抄课件的 $V$，两个含义必须分句说明。}

\begin{center}
\small
\tabtop
\begin{tabular}{@{}p{2.4cm}p{8.0cm}p{4.4cm}@{}}
\thead{符号} & \thead{含义} & \thead{单位} \\
\tabmid
$E$, $V$, $d$ & 外加电场、电压、电极间距（横向几何中 $E=V/d$） & V/m, V, m \\
\tabmid
$L$（或 $l$） & 光在晶体中走过的长度 & m \\
\tabmid
$n$, $n_o$, $n_e$ & 未加场时的折射率；寻常光、非常光折射率 & 无量纲 \\
\tabmid
$\Delta n$ & 双折射（两种正交偏振的折射率差） & 无量纲 \\
\tabmid
$r$（$r_{63}$, $r_{33}$\dots） & 线性电光系数（Pockels 系数） & m/V（常用 pm/V） \\
\tabmid
$K$ & 二次电光系数（Kerr 常数） & m/V$^2$ \\
\tabmid
$\Gamma$（课件也写 $\Delta\phi$） & 相位延迟（retardation） & rad \\
\tabmid
$V_\pi$ & 半波电压：产生 $\Gamma=\pi$ 所需的电压 & V \\
\tabmid
$T$, $I$ & 透射率（归一化强度）、光强 & 无量纲, W/m$^2$ \\
\tabmid
$\vdt$ & Verdet 常数（本讲义专用符号） & rad/(T$\cdot$m) \\
\tabmid
$B$, $M$, $H$ & 磁感应强度、磁化强度、磁场强度 & T, A/m, A/m \\
\tabmid
$\theta$, $\psi$ & 偏振面转角、椭偏率（椭圆率角） & rad \\
\tabmid
$\varepsilon$, $\varepsilon_{xy}$, $Q$ & 介电常数、介电张量的非对角元、磁光常数（$Q=g/\varepsilon$） & 无量纲（相对值） \\
\tabmid
$r_{pp}$, $r_{ps}$ & 反射系数：$p\to p$、$p\to s$（MOKE 用） & 无量纲 \\
\tabmid
$d$, $A$, $g$ & CD 中的光程、吸光度、不对称因子 $g=\Delta A/A$ & m, 无量纲, 无量纲 \\
\end{tabular}
\tabbot
\normalsize
\end{center}

\subhdr{0.3 怎么读这份讲义}

\begin{itemize}
\item 想先抓住骨架：只看 §0 的图和 §1、§2 的判据，再跳到 §12（互易性）——那是全讲最容易考、
      也最容易在作业里翻车的一处。
\item 要写作业：题干在 §16；动笔前先读 §5--§6 的 $\Gamma$ 与 $V_\pi$ 两式（作业第 1 题）、
      §11--§12 的 $\theta=\vdt BL$ 与往返结论（第 2 题）、§2 的对称性判据与 §4 的 $\Delta n=K\lambda E^2$（第 3 题）。
      详解在附录 C，英文答题版在附录 D。
\item 只想过一遍工程应用：§7（调制器、Q 开关）与 §15（隔离器、环形器、磁光存储）。
\end{itemize}

%=====================================================================
\hdr{§1}{外场改的是折射率：三个量、一条换算}

\subhdr{1.1 先看这张分类表（课件 p4）}

课件开篇就把全讲的五种效应按「谁在推、推出来的响应是几次方」摆成一张表。这张表是整讲的目录。

\begin{center}
\small
\tabtop
\begin{tabular}{@{}p{4.0cm}p{4.2cm}p{7.0cm}@{}}
\thead{外部影响} & \thead{效应} & \thead{场的依赖} \\
\tabmid
电场 $E$ & Kerr 效应 & 二次（$E^2$），所有材料 \\
\tabmid
电场 $E$ & Pockels 效应 & 线性（$E$），仅非中心对称 \\
\tabmid
磁场 $B$ & Faraday 效应 & 线性（$B$），透射 \\
\tabmid
磁场 $B$ & Cotton--Mouton（Voigt）效应 & 二次（$B^2$），横向 \\
\tabmid
无外场 & 自然旋光 & 介质自身的结构（手性），与场无关 \\
\end{tabular}
\tabbot
\normalsize
\end{center}

这张表有一个立即可用的推论：\textbf{「线性」总是被对称性限制，「二次」从不被限制}。
原因在 §2，一句话：有反演中心的介质，$\pm E$ 必须给出同一个结果，所以奇次项只能为零。
读表时先把这条记在手边，后面每一节都会用到。

\subhdr{1.2 为什么值得管这件事（课件 p3）}

课件列的动机可以归成三类，正好对应三种「精度要求」：

\begin{itemize}
\item \textbf{要快}：通信里的电光调制器把数据编码到激光上，速率几十 GHz。这要求响应是电子的（fs 级），
      而不是分子转动（ns 级）——§3.2 会看到这个差别把 Kerr 和 Pockels 分到两个应用领域。
\item \textbf{要猛}：Q 开关把腔内储能一次性放出，得到巨脉冲，用于加工与聚变诊断（§7.7）。
\item \textbf{要准}：磁场/电流传感器、陀螺、磁光存储读出头，靠的是「转角与磁场成正比」这一个线性关系（§11、§15）。
\end{itemize}

共同点不是「效应很神奇」，而是\textbf{它们都把「一个电学量」翻译成「一个可以被偏振片读出的角度或相位」}。
整讲的工程内容都是这句话的展开。

\subhdr{1.3 偏振态与琼斯矢量（课件 p5，回顾 L1--L2）}

只回顾本讲要用的三件事。

\textbf{（i）任意线偏振可以拆成一左一右两个圆偏振。}这是 Fresnel 解释旋光的全部基础（§9），
也是磁光效应「把 RCP 与 LCP 的折射率分开」这句话能成立的前提：

\begin{equation}
\hat{e}_x=\frac{1}{\sqrt2}\left[\hat{e}_R+\hat{e}_L\right],\qquad
\hat{e}_R=\frac{1}{\sqrt2}\binom{1}{-i},\qquad
\hat{e}_L=\frac{1}{\sqrt2}\binom{1}{i}.
\end{equation}

\textbf{（ii）两个正交分量之间只差一个相位，合起来就是椭圆。}
写 $\vec E=(E_x,\,E_y e^{i\delta})$，$\delta=0$ 或 $\pi$ 是线偏振，$\delta=\pm\pi/2$ 且 $|E_x|=|E_y|$ 是圆偏振，
其余是椭圆。本讲所有「相位延迟 $\Gamma$」都只是这个 $\delta$ 被改了一个量。

\textbf{（iii）手性约定各教材不同。}L1--L3 已经统一过：本课程按「迎着光看，电矢量逆时针转」记为左旋。
写作业时把约定写在旁边，结论（「等幅 + $90^\circ$ = 圆」）与约定无关。

\subhdr{1.4 双折射与指标体（课件 p6）}

\textbf{折射率}定义相速度：$n=c/v$。\textbf{双折射}是两个正交偏振分量的折射率不同：
$\Delta n=n_{\text{慢}}-n_{\text{快}}$。

为什么用「指标体」（optical indicatrix）说话：折射率是各向异性的，方向不同 $n$ 不同，
一个标量写不下。指标体是一个椭球，
\begin{equation}
\frac{x^2}{n_x^2}+\frac{y^2}{n_y^2}+\frac{z^2}{n_z^2}=1,
\end{equation}
它的三个半轴长度就是三个主轴方向的折射率。波沿某个方向传播时，允许的两个正交偏振方向
就是与该方向垂直的截面椭圆的\emph{两个半轴}方向，相应的半轴长度就是这两个偏振看到的折射率。

\begin{Callout}
\textbf{本讲的统一语言。}电光效应与磁光效应的差别，可以精确地说成「它们各自把指标体怎么改」：
\begin{itemize}
\item Pockels / Kerr：电场把椭球\emph{变形}（轴长变、主轴方向也可能转）——改的是\textbf{线}双折射，
      结果是两个\emph{线}偏振之间出现相位差 $\Gamma$。
\item Faraday：磁场不改变线偏振的两个主轴，而是让右旋与左旋圆偏振的折射率不同——改的是\textbf{圆}双折射，
      结果是整个偏振面\emph{转}一个角 $\theta$。
\end{itemize}
一个是「两个线偏振错开相位」，一个是「两个圆偏振错开相位」——这两个说法在 §12 会合，变成互易与非互易的判据。
\end{Callout}

\subhdr{1.5 从 $\Delta n$ 到可测量的相位差}

这是全讲出现频率最高的一个换算，务必写熟：
\begin{equation}
\label{eq:gamma}
\boxed{\;\Gamma=\frac{2\pi}{\lambda}\,\Delta n\,L\;}
\end{equation}
其中 $\lambda$ 是\emph{真空}波长，$L$ 是光在介质中走过的长度。

推导只有一句：两个分量走过同一段几何路程 $L$，但相速度不同，所以相位各自累积
$\phi_i=2\pi n_iL/\lambda$，相减就得到 (\ref{eq:gamma})。

\begin{WarnBox}
\textbf{三个常见错法。}
（1）把 $\lambda$ 代成介质中的波长 $\lambda/n$——那样 $\Delta n$ 就会被多除一次 $n$；
（2）把 $L$ 代成「样品厚度」但不看光是否斜入射——本讲全部按垂直入射、$L$ 就是电极间那段长度；
（3）把 $\Gamma$ 与「延迟量」$\Delta=\Delta n\,L$（单位是米）混用。$\Delta$ 是光程差，$\Gamma$ 是相位差，两者差一个 $2\pi/\lambda$。
\end{WarnBox}

\textbf{小例（量级感）}：设 $\Delta n=1.0\times10^{-5}$，$L=1\ \mathrm{cm}$，$\lambda=633\ \mathrm{nm}$。

\begin{equation}
\Gamma=\frac{2\pi}{633\times10^{-9}}\times1.0\times10^{-5}\times0.01
=\frac{2\pi\times1.0\times10^{-7}}{6.33\times10^{-7}}\approx0.99\ \mathrm{rad}\approx57^\circ .
\end{equation}

也就是说：\textbf{折射率只要被改动十万分之一，一厘米的晶体就能造出接近 $60^\circ$ 的相位差}。
反过来，$\Gamma$ 可以被检偏器读到 $\lambda/1000$ 的精度，所以 $\Delta n$ 的灵敏度到 $10^{-8}$ 量级——
这就是为什么电光与磁光效应虽然本身很弱，却仍然是精密的测量手段。

%=====================================================================
\hdr{§2}{对称性：为什么有的线性、有的二次}

\subhdr{2.1 判据（课件 p7、p12、p24）}

\begin{Callout}
\textbf{判据。}把折射率的变化写成场的幂级数
\begin{equation}
\Delta n(E)=a_1E+a_2E^2+a_3E^3+\cdots
\end{equation}
若介质具有\textbf{反演中心}（centrosymmetric，即存在一点，把全部坐标 $\vec r\to-\vec r$ 后晶体与自身重合），
则 $a_1=a_3=a_5=\cdots=0$：\textbf{只允许偶次项}。这时只剩 Kerr 效应。

若介质\textbf{没有反演中心}，一次项允许存在，且通常远大于二次项：\textbf{Pockels 效应}。
\end{Callout}

这个判据的力量在于它不需要知道具体的原子结构：只要告诉我晶体有没有反演中心，我就能说出它能不能做普克尔斯盒。

\subhdr{2.2 为什么会这样：完整论证}

\textbf{第一步，物理要求。}反演中心意味着「把晶体绕中心转 $180^\circ$（等价于 $\vec r\to-\vec r$）后，
晶体看上去一模一样」。所以「加 $+E$ 得到的折射率变化」与「加 $-E$ 得到的折射率变化」，
在这块晶体里必须是\emph{同一个物理状态}：电场方向反转了，但晶体也跟着「反转」了，两者互相抵消。

形式化地写：对中心对称介质，
\begin{equation}
\label{eq:inversion}
\Delta n(+E)=\Delta n(-E)\qquad\text{对任意 }E .
\end{equation}

\textbf{第二步，代入级数。}把级数代进 (\ref{eq:inversion})：
\begin{equation}
a_1E+a_2E^2+a_3E^3+\cdots=a_1(-E)+a_2E^2+a_3(-E)^3+\cdots
\end{equation}
两边相减，偶次项全部消掉，剩下
\begin{equation}
2a_1E+2a_3E^3+2a_5E^5+\cdots=0\qquad\text{对任意 }E .
\end{equation}

\textbf{第三步，结论。}一个多项式对\emph{任意} $E$ 恒为零，只能所有系数为零：
$a_1=a_3=a_5=\cdots=0$。若 $a_1\neq0$，就要求 $2a_1E\equiv0$，与「任意 $E$」矛盾。
所以线性项必被禁掉，最低阶允许项是 $a_2E^2$——这就是 Kerr 效应。

\aside{注意 (\ref{eq:inversion}) 里没有写成「$\Delta n$ 的符号翻转」。$n$ 是正数，$n$ 的变化量可以正可以负，
但「$+E$ 与 $-E$ 给出同一个变化量」这句话是实的；Pockels 效应的特征恰恰是 $\Delta n(-E)=-\Delta n(+E)$
（符号跟着场翻），这正是它被中心对称禁掉的原因。}

\begin{center}
\includegraphics[width=0.62\linewidth]{fig-l4/fig1-dn-vs-E.png}

{\small\color{inkgray}图 0：奇函数与偶函数。两条曲线各自归一化到 $E=+E_0$ 处的值，所以只比较\emph{形状}：
Pockels（深紫）关于原点反对称，反接电压就变色；Kerr（浅紫）关于竖轴对称，反接电压毫无变化。
两者在实际量级上差约 $10^3$ 倍（§5.5），这里为看清对称性而做了归一化。}
\end{center}

\subhdr{2.3 小例：若真有线性项，会被看见吗}

设某介质的形式系数 $a_1=1.0\times10^{-12}\ \mathrm{m/V}$、$a_2=1.0\times10^{-20}\ \mathrm{m/V^2}$，
在 $E=1.0\times10^{6}\ \mathrm{V/m}$ 处：

\begin{equation}
a_1E=1.0\times10^{-6},\qquad a_2E^2=1.0\times10^{-8}.
\end{equation}

线性项比二次项大 $100$ 倍。假如它存在于一块有反演中心的晶体里，那么把电压反接，
$\Delta n$ 会从 $+1.0\times10^{-6}$ 变成 $-1.0\times10^{-6}$——这在偏光显微镜下是肉眼可见的差别。
既然实验上观察不到这种极性依赖，唯一的解释就是 $a_1$ 确实为零。这个数字对比也解释了另一件事：
\textbf{一旦允许线性项，二次项就被压得可以忽略}（§3.1 会把它写成 $10^3$ 倍的量级差）。

\subhdr{2.4 材料类型表（课件 p8）}

\begin{center}
\small
\tabtop
\begin{tabular}{@{}>{\raggedright}p{7.4cm}>{\raggedright}p{3.8cm}>{\raggedright\arraybackslash}p{4.2cm}@{}}
\thead{材料类型} & \thead{Pockels（线性）} & \thead{Kerr（二次）} \\
\tabmid
中心对称（玻璃、液体、Si、金刚石、NaCl） & 禁止 & 允许 \\
\tabmid
非中心对称（KDP、LiNbO$_3$、石英、GaAs） & 允许 & 也同时存在 \\
\end{tabular}
\tabbot
\normalsize
\end{center}

\subhdr{2.5 判据的两个边界（别用过头）}

第一，\textbf{「非中心对称」是必要条件，不是充分条件}。KDP 属于 $\bar42m$ 点群，有 $r_{63}$；
但同一族里有的切型其有效系数为零，必须看张量的哪个分量被用上（§5.2）。

第二，\textbf{课件特别点出 GaAs 是个「例外」}：它属于闪锌矿结构（$\bar43m$），没有反演中心，
所以既有 Pockels 也有 Kerr。课件 p8 的脚注说它「重要」，是因为它在半导体激光与红外调制里常用；
写「GaAs 是非中心对称的」这句话时要分清它\emph{看起来}像立方对称、但点群确实无中心。

\begin{WarnBox}
\textbf{一句话记住。}中心对称 $\Rightarrow$ 只能二次；非中心对称 $\Rightarrow$ 线性可用，二次仍存在。
判断一块材料能不能做普克尔斯盒，只看\emph{点群}，不看硬度、透明范围、折射率大小——那些是第二步的工程选择（§6.3）。
\end{WarnBox}

%=====================================================================
\hdr{§3}{Kerr 效应：折射率随 $E^2$ 变}

\subhdr{3.1 公式与量纲（课件 p10）}

\begin{equation}
\label{eq:kerr}
\boxed{\;\Delta n=K\,\lambda\,E^2\;}
\end{equation}
其中 $K$ 是 Kerr 常数（材料给定，单位 $\mathrm{m/V^2}$），$\lambda$ 是真空波长，$E$ 是外加场强。

量纲先查一遍，免得把 $K$ 的单位记错：
\begin{equation}
[K]\,[\lambda]\,[E^2]=\frac{\mathrm{m}}{\mathrm{V^2}}\cdot\mathrm{m}\cdot\frac{\mathrm{V^2}}{\mathrm{m^2}}=1 .
\end{equation}
三个量的单位必须这样配平，结果才是无量纲的折射率差。这一步在作业第 3 题里会直接用到。

\textbf{两条立刻能用的推论。}

（i）\textbf{反接电压不变}。$E\to-E$ 时 $\Delta n$ 不变，因为 $(-E)^2=E^2$。
换成 AC 驱动，$\Delta n$ 会以\emph{两倍}驱动频率振荡——这是 Kerr 盒做高压测量时的指纹（§4.4）。

（ii）\textbf{两次方意味着「要么很长的 $L$、要么很高的 $E$」}。把 (\ref{eq:kerr}) 代进 (\ref{eq:gamma})：
\begin{equation}
\label{eq:gamma-kerr}
\Gamma=\frac{2\pi}{\lambda}\,\Delta n\,L=2\pi K L E^2=2\pi K L\left(\frac{V}{d}\right)^{2}.
\end{equation}
注意 $\lambda$ \emph{消掉了}——Kerr 盒的延迟与光的波长无关（在 $K$ 本身随波长缓变的意义上）。
这一点与 Pockels 的 $V_\pi\propto\lambda$ 很不一样，写作业第 3 题时值得顺手提一句。

\subhdr{3.2 两种机理与两个时间尺度（课件 p11）}

同一个 $E^2$ 现象背后有两种完全不同的物理，它们的时间尺度差三个数量级：

\begin{center}
\small
\tabtop
\begin{tabular}{@{}p{3.0cm}p{5.4cm}p{6.2cm}@{}}
\thead{机理} & \thead{在做什么} & \thead{响应时间} \\
\tabmid
分子取向 & 各向异性分子有永久偶极矩，场把它「拧」到同一个方向，宏观上出现各向异性 & 纳秒（由转动扩散限制；$\sim10\ \mathrm{ns}$ for 硝基苯） \\
\tabmid
电子畸变 & 场把原子/分子的电子云拉变形，诱导出偶极矩 & 飞秒（电子极化） \\
\end{tabular}
\tabbot
\normalsize
\end{center}

\textbf{哪个占主导？}液体里是取向（分子能动，场能把它排队），固体里是电子（分子被晶格钉住，只能变形）。
由此得到整讲反复用到的一条对应：
\begin{equation}
\text{液体：}K\text{ 大（}\sim10^{-12}\ \mathrm{m/V^2}\text{）、慢（ns）};\qquad
\text{固体：}K\text{ 小（}\sim10^{-16}\ \mathrm{m/V^2}\text{）、快（fs）}.
\end{equation}

\subhdr{3.3 Kerr 常数的实际量级（课件 p13）}

\begin{center}
\small
\tabtop
\begin{tabular}{@{}p{5.0cm}p{3.6cm}p{6.0cm}@{}}
\thead{材料} & \thead{$K$（$\mathrm{m/V^2}$）} & \thead{备注} \\
\tabmid
硝基苯 C$_6$H$_5$NO$_2$ & $\sim2.2\times10^{-12}$ & 已知最大；标准 Kerr 盒液体 \\
\tabmid
二硫化碳 CS$_2$ & $\sim3.6\times10^{-14}$ & 常用但有毒 \\
\tabmid
水 H$_2$O & $\sim4.4\times10^{-14}$ & 很弱；对生物重要 \\
\tabmid
苯 C$_6$H$_6$ & $\sim0.7\times10^{-14}$ & 中等 \\
\tabmid
火石玻璃 & $\sim10^{-16}$ & 很弱；需高场 \\
\tabmid
空气（1 atm） & $\sim10^{-19}$ & 极弱 \\
\end{tabular}
\tabbot
\normalsize
\end{center}

读这张表要看出\textbf{四个数量级}的跨度：硝基苯比玻璃大约 $10^4$ 倍。这个比值直接决定了 §4.1 的工程结论：
液体做 Kerr 盒是「大效应 + 慢 + 麻烦」，固体做 Kerr 盒是「快 + 弱到几乎不能用」。

\subhdr{3.4 温度与波长依赖（课件 p14）}

\textbf{温度}：液体里 $K$ 随温度下降，因为热运动在破坏分子的定向排列，近似有
\begin{equation}
K\propto\frac{1}{T}\qquad\text{（取向机理主导时）}.
\end{equation}
工程含义很硬：硝基苯在 $\sim5.7^\circ\mathrm{C}$ 就冻住了，这就是这类 Kerr 盒的\emph{工作下限}。

\textbf{波长}：$K$ 一般随波长变短而增大，经验上 $K\propto\lambda^{-2}$ 到 $\lambda^{-3}$（随材料不同），
根源是折射率色散。把这一条和 (\ref{eq:kerr}) 里的那个 $\lambda$ 合起来看：
$\Delta n=K(\lambda)\lambda E^2$，所以 $\Delta n$ 对波长的依赖弱于「正比于 $\lambda$」，
但仍不是完全相同——课件把 $\lambda$ 显式写在公式里，是为了让 $K$ 的单位与「不同波长可比」这两件事都成立。

\subhdr{3.5 为什么 Kerr 效应在任何材料里都存在}

回到 §2 的判据：$E^2$ 是偶函数，满足 (\ref{eq:inversion})，所以中心对称介质不禁它。这就解释了课件 p12 的清单：
液体（硝基苯、CS$_2$、水）、气体（强场下）、玻璃与无定形体、中心对称晶体（Si、金刚石、NaCl）全都有 Kerr 效应。

\textbf{反过来说}：对于\textbf{中心对称}材料，Kerr 是\emph{唯一}可用的电光效应。这句话在 §4.4 的第 2 问里会变成
「既然 Pockels 更强，为什么还留着 Kerr 盒」的答案。

%=====================================================================
\hdr{§4}{Kerr 盒：数值实感、几何与三问}

\subhdr{4.1 液体盒与固体盒（课件 p15）}

\begin{center}
\small
\tabtop
\begin{tabular}{@{}p{4.0cm}p{5.2cm}p{5.2cm}@{}}
\thead{性质} & \thead{液体（硝基苯）} & \thead{固体（玻璃）} \\
\tabmid
Kerr 常数 & 大（$10^{-12}\ \mathrm{m/V^2}$） & 小（$10^{-16}\ \mathrm{m/V^2}$） \\
\tabmid
响应时间 & 纳秒（取向） & 飞秒（电子） \\
\tabmid
所需电压 & 中等（kV 级） & 极高（MV 级） \\
\tabmid
实用盒长 & 厘米 & 毫米 \\
\tabmid
维护 & 需密封 & 结实、紧凑 \\
\tabmid
温度稳定 & 差（冻结、挥发） & 极好 \\
\end{tabular}
\tabbot
\normalsize
\end{center}

\subhdr{4.2 盒的几何（课件 p16）}

标准 Kerr 盒用\textbf{横向几何}：电场垂直于光的传播方向，两块平行板电极浸在液体里；
孔径受电极间距限制（典型 $1$--$10\ \mathrm{mm}$），通光长度 $1$--$10\ \mathrm{cm}$ 才够累出可用的延迟；
前后各放一片偏振器、且相互正交，用来把相位差转成强度变化。

\begin{center}
\begin{tikzpicture}[font=\footnotesize,x=1cm,y=1cm]
% 光轴
\draw[->,inkgray] (-0.4,0) -- (14.2,0) node[right,text=inkgray]{\scriptsize 光轴};
% 入射光（非偏振）
\foreach \a in {20,55,90,125,160}{\draw[->,red!65!black] (0,0) -- ++(\a:0.42);}
\node[align=center] at (0.1,-0.95) {\scriptsize 非偏振光};
% 起偏器
\draw[fill=bandcream,draw=petal] (1.6,-0.85) rectangle (1.95,0.85);
\node[align=center] at (1.78,-1.25) {\scriptsize 起偏器};
% 偏振光箭头（竖直）
\draw[<->,violet,thick] (2.7,0) -- (2.7,0.55);
\draw[<->,violet,thick] (3.2,0) -- (3.2,0.55);
% Kerr 盒
\draw[draw=violet,thick] (4.0,-1.0) rectangle (9.4,1.0);
\draw[fill=bandviolet,draw=violet] (4.0,1.0) rectangle (9.4,1.22);
\draw[fill=bandviolet,draw=violet] (4.0,-1.22) rectangle (9.4,-1.0);
\node at (6.7,1.11) {\scriptsize $+$};
\node at (6.7,-1.11) {\scriptsize $-$};
% 场线
\foreach \x in {4.6,5.4,6.2,7.0,7.8,8.6}{
  \draw[->,violet] (\x,1.0) -- (\x,-1.0);}
\node[align=center,text=violet] at (6.7,-1.75) {\scriptsize $E=V/d$，垂直于光传播方向};
\node[align=center] at (6.7,1.62) {\scriptsize 电极（平行板）};
% 出射光（椭圆）
\draw[->,violet,thick] (10.0,0) .. controls (10.25,0.34) and (10.5,-0.34) .. (10.75,0.02);
\draw[->,violet,thick] (10.75,0.02) .. controls (11.0,0.38) and (11.25,-0.30) .. (11.5,0.03);
\node[align=center] at (10.85,-0.95) {\scriptsize 椭圆偏振};
% 检偏器（与起偏器正交）
\draw[fill=bandcream,draw=petal] (12.1,-0.85) rectangle (12.45,0.85);
\draw[<->,violet,thick] (12.0,-0.5) -- (12.55,0.5);
\node[align=center] at (12.28,-1.25) {\scriptsize 检偏器（正交）};
\draw[->,red!65!black] (13.2,0) -- (14.0,0) node[below,text=black]{\scriptsize 探测器};
\end{tikzpicture}
\end{center}

{\small\color{inkgray}图 1：Kerr 盒的横向几何。电场在电极板之间竖直向下，光沿水平方向穿过，
两者垂直；因此 $E=V/d$ 里的 $d$ 是\emph{电极间距}，不是盒长。盒长 $L$ 决定累积多少延迟。}

\subhdr{4.3 响应时间与速度上限（课件 p17）}

\begin{center}
\small
\tabtop
\begin{tabular}{@{}p{6.0cm}p{4.2cm}p{4.6cm}@{}}
\thead{器件} & \thead{响应时间} & \thead{瓶颈} \\
\tabmid
液体 Kerr 盒 & $1$--$100\ \mathrm{ns}$ & 分子转动扩散（黏度、温度） \\
\tabmid
固体 Kerr 盒 & $<1\ \mathrm{ps}$ & 电子极化（极快但需极高场） \\
\tabmid
Pockels 盒 & ps--ns & 电子 + 声子 \\
\end{tabular}
\tabbot
\normalsize
\end{center}

结论一句话：\textbf{要 GHz 以上调制，Pockels 盒占主导；Kerr 盒守着自己那块「大效应、大孔径、不挑极性」的地盘}（§4.4 第 2 问）。

\subhdr{4.4 课件三问（p18--p20）}

\textbf{问 1：既然是二次方，反接电压是不是没有影响？}是。$(-E)^2=E^2$，$\Delta n$ 不变。
实践后果：用 AC 驱动 Kerr 盒时，\emph{正负半周给出同样的光学效果}，所以 $T(t)$ 以两倍频振荡。
这与 Pockels 盒相反——后者反接电压就反转 $\Delta n$ 的符号。

\textbf{问 2：Kerr 弱得多，为什么还留着它？}两个理由。其一，$E^2$ 是二阶过程，同场强下比 $E$ 小几个数量级；
液体里又是「热运动对抗定向」的效率不高的机理。其二，它仍然不可替代的场合有四类：
必须用中心对称材料时；测极高电压时（二次响应便于线性化读数）；需要大孔径液体介质时；
以及超快固体 Kerr 门（光学 Kerr 效应，用于飞秒脉冲测量）。

\textbf{问 3：光能不能沿着电场走？}标准盒子里不能——横向几何下，场\emph{定义}了一个光轴：
平行于 $E$ 的偏振与垂直于 $E$ 的偏振看到不同折射率，这才是双折射的来源。若光平行于 $E$，
它在横向平面内看到的是各向同性的（两个垂直方向等价），双折射为零。这条与 §6.1 的「纵向/横向」是同一件事的电场版。

\subhdr{4.5 小例：一个硝基苯 Kerr 盒要多少电压}

设液体为硝基苯，$K=2.2\times10^{-12}\ \mathrm{m/V^2}$，$L=1\ \mathrm{cm}$，极板间距 $d=2\ \mathrm{mm}$，
要求四分之一波延迟 $\Gamma=\pi/2$。由 (\ref{eq:gamma-kerr}) 反解：
\begin{equation}
\left(\frac{V}{d}\right)^2=\frac{\Gamma}{2\pi KL}
\;\Longrightarrow\;
V=d\sqrt{\frac{\Gamma}{2\pi K L}} .
\end{equation}
先算分母：$2\pi KL=2\pi\times2.2\times10^{-12}\times0.01\approx1.38\times10^{-13}\ \mathrm{m/V^2}$，
\begin{equation}
\frac{V}{d}=\sqrt{\frac{\pi/2}{1.38\times10^{-13}}}\approx\sqrt{1.14\times10^{13}}
\approx3.4\times10^{6}\ \mathrm{V/m},
\qquad
V\approx3.4\times10^{6}\times2\times10^{-3}\approx6.8\times10^{3}\ \mathrm{V}.
\end{equation}

\textbf{约 $6.7\ \mathrm{kV}$}，这才是硝基苯盒子能工作的量级。把同一个式子用在弱液体
（$K\sim2\times10^{-14}$，即比硝基苯小 $100$ 倍）上，同样的 $d=2\ \mathrm{mm}$、$L=1\ \mathrm{cm}$ 就要 $7.1\times10^{4}\ \mathrm{V}$；
再把间隙放大到 $d=1\ \mathrm{cm}$、盒子加长到 $L=5\ \mathrm{cm}$，就是 $1.5\times10^{5}\ \mathrm{V}$ 量级——
课件 p46--p48 的例题正是这一组（见 §8.2）。

两条经验因此成立：\textbf{（i）$V\propto d$，把电极间隙做小是降电压最有效的手段；
（ii）$V\propto1/\sqrt{L}$，加长盒子只按平方根降电压}——与 Pockels 的 $V_\pi\propto d/L$（线性）不同（§6.2）。

%=====================================================================
\hdr{§5}{Pockels 效应：折射率随 $E$ 线性变}

\subhdr{5.1 公式从哪里来（课件 p22）}

\begin{equation}
\label{eq:pockels}
\boxed{\;\Delta n=-\frac{1}{2}n^3rE\;}
\end{equation}

这个式子里的 $\tfrac12$ 和 $n^3$ 不是约定，而是从「指标体」的定义算出来的。设电场让
\emph{折射率平方的倒数}改变一个与场成正比的小量：
\begin{equation}
\label{eq:invn}
\frac{1}{n'^2}=\frac{1}{n^2}+rE .
\end{equation}
这是电光效应最自然的写法，因为指标体方程里出现的正是 $1/n^2$（§1.4）。
现在令 $n'=n+\delta n$，把左边展开到一阶：
\begin{equation}
\frac{1}{(n+\delta n)^2}=\frac{1}{n^2}\left(1+\frac{\delta n}{n}\right)^{-2}
\approx\frac{1}{n^2}\left(1-2\frac{\delta n}{n}\right)
=\frac{1}{n^2}-\frac{2\,\delta n}{n^3} .
\end{equation}
与 (\ref{eq:invn}) 比较，$1/n^2$ 抵消，剩下
\begin{equation}
-\frac{2\,\delta n}{n^3}=rE
\;\Longrightarrow\;
\delta n=-\frac{1}{2}n^3rE .
\end{equation}

所以：\textbf{$n^3$ 来自对 $1/n^2$ 求导，$\tfrac12$ 来自二次项的泰勒展开}。
记住这个来历，就不会在作业里把 $n^3$ 写成 $n$、或把 $\tfrac12$ 丢掉。

\begin{WarnBox}
\textbf{最容易错的一处：$\tfrac12$ 在算 $\Gamma$ 时会「消失」。}
式 (\ref{eq:pockels}) 说的是\emph{一个}偏振分量的折射率变化量。可是横向 Pockels 盒里，
两个正交偏振看到的是 $n+\tfrac12n^3rE$ 与 $n-\tfrac12n^3rE$（符号相反），
所以\textbf{两者之差}是
\begin{equation}
\label{eq:dnalpha}
\Delta n_{\text{两模之间}}=n^3rE ,
\end{equation}
$\tfrac12$ 被两个模的差抵消掉了。把 (\ref{eq:dnalpha}) 代进 (\ref{eq:gamma}) 才得到课件的
$\Gamma=\frac{2\pi}{\lambda}n^3rV\frac{L}{d}$。若误用 (\ref{eq:pockels}) 直接当 $\Delta n$，结果会差一倍——
这是本章最常见的失分点。
\end{WarnBox}

\subhdr{5.2 张量形式：为什么必须写 $r_{ij}$（课件 p25）}

Pockels 效应是各向异性的：给定电场方向，指标体会沿特定方向变形，变形量取决于晶体的取向。
严格的写法是对\emph{逆介电张量}的六个独立分量逐项线性响应：
\begin{equation}
\Delta\!\left(\frac{1}{n^2}\right)_{\!i}=\sum_{j=1}^{3}r_{ij}E_j,
\qquad i=1,\dots,6;\ j=x,y,z,
\end{equation}
$r_{ij}$ 是一个 $6\times3$ 的矩阵（原本是三阶张量 $r_{ijk}$ 的 27 个分量，用 Voigt 记法压成 18 个），
其非零元由晶体的点群决定。

\textbf{KDP 的情形}（$\bar42m$ 点群）：只有 $r_{41}=r_{52}$ 与 $r_{63}$ 非零。做纵向调制器时用的是 $r_{63}$：
电场沿 $z$（光轴），光也沿 $z$ 传播，感生的主轴在 $xy$ 平面内\emph{转过 $45^\circ$}——
原来沿 $x,y$ 的两个主轴变成沿 $x\pm y$ 方向。这件事的物理含义是：
\textbf{一个各向同性（或单轴）的晶体被场变成了双轴晶体}，而双折射的大小正比于 $E$。

\begin{Callout}
\textbf{为什么「转 $45^\circ$」很要紧。}如果入射的线偏振沿 $x$（原来的主轴），它在这个新坐标系里
是 $45^\circ$ 方向的等幅叠加，于是被完整地「劈」成两束等幅分量（$\cos45^\circ=\sin45^\circ$），
延迟 $\Gamma$ 才能在后面的检偏器上全部转化为强度调制。若偏振沿新主轴，则只看到一个折射率，
延迟为零，调制器完全不工作。§7.1 的装置图里那个「$45^\circ$」就是这么来的。
\end{Callout}

\subhdr{5.3 材料与对称类（课件 p23）}

\begin{center}
\small
\tabtop
\begin{tabular}{@{}p{4.2cm}p{4.4cm}p{5.8cm}@{}}
\thead{晶类（点群）} & \thead{例子} & \thead{电光系数} \\
\tabmid
四方 $\bar42m$ & KDP（KH$_2$PO$_4$） & $r_{63}\approx10.5\ \mathrm{pm/V}$ \\
\tabmid
四方 $\bar42m$ & KD$^*$P（氘化） & $r_{63}\approx26\ \mathrm{pm/V}$ \\
\tabmid
三方 $3m$ & LiNbO$_3$ & $r_{33}\approx30.8\ \mathrm{pm/V}$ \\
\tabmid
立方 $\bar43m$ & GaAs & $r_{41}\approx1.4\ \mathrm{pm/V}$ \\
\tabmid
三方 $32$ & 石英 & $r_{11}\approx0.2\ \mathrm{pm/V}$ \\
\end{tabular}
\tabbot
\normalsize
\end{center}

课件在这张表下面写了一句「规则」：\textbf{有反演中心，Pockels 效应严格为零}。
中心对称材料（Si、Ge、NaCl、金刚石、玻璃、液体、气体）只有 Kerr 效应可用。

表里 $r$ 的单位是 $\mathrm{pm/V}=10^{-12}\ \mathrm{m/V}$。KDP 的 $10.5$ 与 §6.3 表里的 $10.5$ 是同一个数（课本常写 $10.6$）。

\subhdr{5.4 课件 FAQ：原子层面到底发生了什么（p33）}

问：为什么中心对称晶体里线性响应被禁，原子尺度上「禁」在哪里？

答分两层。\textbf{数学层}在 §2.2 已经证完：$\Delta n(+E)=\Delta n(-E)$ 强制 $a_1=0$。
\textbf{原子层}是这件事的实现方式：在非中心对称的晶体里，电子云与离子的排布本身就是不对称的，
场一推，电子分布沿场方向平移、离子沿同一方向位移，两者的位移都\emph{正比于}场——于是介电响应出现一项 $\propto E$。
在中心对称的晶体里，场把电子往一边推、把相邻的离子往反方向对称地推，
一阶响应正好相互抵消，剩下的第一项非零贡献只能是 $E^2$。

两层的对应关系值得记：\textbf{「一阶响应相互抵消」正是 (\ref{eq:inversion}) 那条恒等式的物理内容}。
数学上写 $a_1=0$，物理上说「对称的推动给不出净的极化方向」。

\subhdr{5.5 与 Kerr 的量级（课件 p31）}

\begin{center}
\small
\tabtop
\begin{tabular}{@{}p{4.0cm}p{5.0cm}p{4.6cm}@{}}
\thead{性质} & \thead{Pockels（线性）} & \thead{Kerr（二次）} \\
\tabmid
场依赖 & $\Delta n\propto E$ & $\Delta n\propto E^2$ \\
\tabmid
对称性要求 & 仅非中心对称 & 所有材料 \\
\tabmid
典型强度 & 强（$r\sim1$--$30\ \mathrm{pm/V}$） & 弱（$K\sim10^{-14}$--$10^{-12}\ \mathrm{m/V^2}$） \\
\tabmid
响应时间 & ps（电子） & ns（液体）/ fs（固体） \\
\tabmid
驱动电压 & 中（$10^2$--$10^4\ \mathrm{V}$） & 高（$10^4$--$10^6\ \mathrm{V}$） \\
\tabmid
极性 & $\Delta n$ 随极性反号 & 与极性无关 \\
\tabmid
常用材料 & LiNbO$_3$、KDP、GaAs & 硝基苯、CS$_2$、玻璃 \\
\end{tabular}
\tabbot
\normalsize
\end{center}

\textbf{量级对比给个数}：在 $E=10^6\ \mathrm{V/m}$、$n=1.5$、$r=10\ \mathrm{pm/V}$ 下，
Pockels 的 $\Delta n\approx\tfrac12n^3rE\approx1.7\times10^{-5}$；
而 $\Delta n=K\lambda E^2$ 取 $K=2\times10^{-14}$、$\lambda=633\ \mathrm{nm}$ 时是 $1.3\times10^{-8}$。
\textbf{相差约 $10^3$ 倍}——这就是「能用 Pockels 就不用 Kerr」的全部算术依据（也是作业第 3 题的一个自然对照）。

%=====================================================================
\hdr{§6}{纵向与横向、半波电压、材料选择}

\subhdr{6.1 两种几何与两个公式（课件 p26）}

\begin{center}
\begin{tikzpicture}[font=\footnotesize,x=1cm,y=1cm]
% ---- (a) 纵向 ----
\node at (3.2,1.5) {\textbf{(a) 纵向：$\boldsymbol{E}\parallel\boldsymbol{k}$}};
\draw[draw=violet,thick,fill=bandviolet] (1.4,-0.85) rectangle (5.0,0.85);
\draw[draw=violet,thick,fill=petal!40] (1.4,-0.85) rectangle (1.52,0.85);
\draw[draw=violet,thick,fill=petal!40] (4.88,-0.85) rectangle (5.0,0.85);
\draw[->,red!65!black,thick] (0.1,0) -- (5.6,0);
\foreach \x in {1.9,2.6,3.3,4.0,4.7}{\draw[->,violet] (\x,0.6) -- (\x,-0.6);}
\node at (2.05,-1.15) {\scriptsize 透明电极};
\node at (4.4,-1.15) {\scriptsize 透明电极};
\node at (3.2,-1.15) {\scriptsize $L$};
\draw[<->,inkgray] (1.52,-1.6) -- (4.88,-1.6);
\node[align=center] at (3.2,-2.2)
  {\scriptsize $\Gamma=\dfrac{2\pi}{\lambda}n^3rV$\\[2pt]\scriptsize 与 $L$、$d$ 都无关};

% ---- (b) 横向 ----
\node at (11.6,1.5) {\textbf{(b) 横向：$\boldsymbol{E}\perp\boldsymbol{k}$}};
\draw[draw=violet,thick,fill=bandviolet] (9.8,-0.85) rectangle (13.4,0.85);
\draw[draw=violet,thick,fill=petal!40] (9.8,0.85) rectangle (13.4,1.02);
\draw[draw=violet,thick,fill=petal!40] (9.8,-1.02) rectangle (13.4,-0.85);
\draw[->,red!65!black,thick] (8.5,0) -- (14.0,0);
\foreach \x in {10.3,11.0,11.7,12.4,13.1}{\draw[->,violet] (\x,0.85) -- (\x,-0.85);}
\node at (11.6,-1.45) {\scriptsize 电极（不透明，在侧面）};
\draw[<->,inkgray] (13.85,-1.02) -- (13.85,1.02);
\node[align=center] at (14.7,0) {\scriptsize $d$};
\node[align=center] at (11.6,-2.35)
  {\scriptsize $\Gamma=\dfrac{2\pi}{\lambda}n^3rV\dfrac{L}{d}$\\[2pt]\scriptsize 靠 $L/d$ 降电压};
\end{tikzpicture}
\end{center}

{\small\color{inkgray}图 2：两种 Pockels 盒几何。纵向的电极必须透明（ITO 或薄金属），
场与光同向，走过的长度与电极间距是同一个量，所以 $L/d=1$，公式里看不出几何；
横向的电极在两侧，$d$ 是电极间距、$L$ 是通光长度，两者可以独立取，于是出现 $L/d$ 这个「免费增益」。}

两者的推导只差一步：$E$ 的位置。

\textbf{纵向}：$E=V/L$，代入 $\Gamma=\frac{2\pi}{\lambda}(n^3rE)L$：
\begin{equation}
\Gamma_{\parallel}=\frac{2\pi}{\lambda}(n^3r)\frac{V}{L}\,L=\frac{2\pi}{\lambda}n^3rV .
\end{equation}
$L$ 被消掉了——\textbf{纵向的延迟只看电压，与晶体多长无关}。

\textbf{横向}：$E=V/d$，但光走过的还是 $L$：
\begin{equation}
\Gamma_{\perp}=\frac{2\pi}{\lambda}(n^3r)\frac{V}{d}\,L=\frac{2\pi}{\lambda}n^3rV\frac{L}{d} .
\end{equation}

\subhdr{6.2 半波电压 $V_\pi$（课件 p28）}

定义：使 $\Gamma=\pi$ 所需的电压。

\begin{center}
\small
\tabtop
\begin{tabular}{@{}p{4.4cm}p{6.6cm}p{3.6cm}@{}}
\thead{几何} & \thead{$V_\pi$} & \thead{依赖} \\
\tabmid
纵向 & $V_\pi=\dfrac{\lambda}{2n^3r}$ & 只有材料与波长 \\
\tabmid
横向 & $V_\pi=\dfrac{\lambda}{2n^3r}\cdot\dfrac{d}{L}$ & 可用 $d/L$ 压低 \\
\end{tabular}
\tabbot
\normalsize
\end{center}

推导就是把 $\Gamma=\pi$ 代进上一节两式再解 $V$。以横行为例：
\begin{equation}
\pi=\frac{2\pi}{\lambda}n^3rV_\pi\frac{L}{d}
\;\Longrightarrow\;
V_\pi=\frac{\lambda}{2n^3r}\cdot\frac{d}{L} .
\end{equation}

\textbf{一条常用的换算}：任何 Pockels 盒的延迟都可以写成
\begin{equation}
\label{eq:VpiV}
\Gamma=\pi\frac{V}{V_\pi},
\end{equation}
因为 $V_\pi$ 的定义正是「让 $\Gamma$ 到 $\pi$ 的电压」，而 $\Gamma$ 线性正比于 $V$。
这个关系在做作业第 1 题时比直接代公式更快，也自动保证 (a)(b) 两问自洽。

\begin{Callout}
\textbf{$V_\pi$ 是「材料的固有属性」还是「器件的属性」？}纵向的是前者（只看 $\lambda,n,r$），
横向的是后者（还看 $d/L$）。这是选型的第一条：\textbf{要低电压就选横向}，
因为长晶体、窄电极在工程上都能做到。

课件的表还要再补一句提醒：表里的 $V_\pi$ 是\textbf{典型器件值}，不是由同一组 $n,r$ 算出来的。
例如 KDP「$7.5\ \mathrm{kV}$ @ $633\ \mathrm{nm}$」这条，用 $\lambda/(2n^3r)$ 在 $546\ \mathrm{nm}$ 处正好得 $7.55\ \mathrm{kV}$，
在 $633\ \mathrm{nm}$ 处则得 $8.75\ \mathrm{kV}$。做作业一律用公式算，不要照抄表（详见附录 B 第 3 条）。
\end{Callout}

\subhdr{6.3 常用材料（课件 p29）}

\begin{center}
\small
\tabtop
\begin{tabular}{@{}p{3.4cm}p{2.2cm}p{2.4cm}p{2.6cm}p{4.0cm}@{}}
\thead{材料} & \thead{$r$ (pm/V)} & \thead{$V_\pi$ (kV)} & \thead{透光 ($\mu$m)} & \thead{备注} \\
\tabmid
KDP & 10.5 & 7.5 & 0.2--1.5 & 经典；易潮解 \\
\tabmid
KD$^*$P & 26 & 3.5 & 0.2--1.5 & $r$ 更大；Q 开关常用 \\
\tabmid
LiNbO$_3$ & 30.8 & 0.4（横向） & 0.4--4.5 & 工业标准；结实 \\
\tabmid
LiTaO$_3$ & 30.5 & 0.4（横向） & 0.4--4.5 & 与 LiNbO$_3$ 相似 \\
\tabmid
GaAs & 1.4 & 5 & 1--12 & 红外；半导体工艺 \\
\tabmid
BBO & 2.2 & 4 & 0.2--2 & 紫外 \\
\tabmid
石英 & 0.2 & 50 & 0.2--2 & 极弱；做传感器 \\
\end{tabular}
\tabbot
\normalsize
\end{center}

选材料的顺序（工程视角，比背表有用）：\textbf{一、点群允许 Pockels 吗}（§2）→
\textbf{二、工作波长透不透}（这一列决定能不能用）→ \textbf{三、$V_\pi$ 与几何能否做到}（§6.2）→
\textbf{四、稳定性：潮解（KDP）、光折变（LiNbO$_3$ 高功率）、温度漂移}（§7.6）。

\subhdr{6.4 横向为什么能降电压，以及它付出的代价}

把 $V_\pi=\frac{\lambda}{2n^3r}\frac{d}{L}$ 拆开看两个可调量：
\begin{itemize}
\item 增大 $L$：晶体做长（例如 $2\ \mathrm{cm}$），电压按 $1/L$ 降；
\item 减小 $d$：电极间距做窄（波导器件里 $d$ 可到几微米），电压按 $d$ 降。
\end{itemize}
乘积效应很猛：$L/d$ 从 $1$ 变到 $10$，电压就降十倍；波导 LiNbO$_3$ 的 $V_\pi$ 可以做到 $1\ \mathrm{V}$ 量级。

代价有两个。第一，\textbf{温度}：横向盒的延迟建立在晶体的天然双折射上（$\Delta n$ 有温度系数），
偏置点会漂；纵向的天然双折射对延迟没有贡献（光沿光轴走，两个方向同 $n$），所以温度稳定性更好。
第二，\textbf{电极与光路的耦合}：横向要有足够大的通光孔径又不能挡住光，电极设计是主要工艺难点。

\subhdr{6.5 课件 FAQ：纵横怎么选（p34）}

\begin{center}
\small
\tabtop
\begin{tabular}{@{}p{3.4cm}p{5.2cm}p{5.2cm}@{}}
\thead{项目} & \thead{纵向} & \thead{横向} \\
\tabmid
场与光 & 平行 & 垂直 \\
\tabmid
电极 & 透明（ITO、薄金属） & 不透明，放侧面 \\
\tabmid
延迟 & $\Gamma=\frac{2\pi}{\lambda}n^3rV$ & $\Gamma=\frac{2\pi}{\lambda}n^3rV\frac{L}{d}$ \\
\tabmid
$V_\pi$ & 与长度无关 & 可用 $L/d$ 压低 \\
\tabmid
孔径 & 受电极限制 & 可做大 \\
\tabmid
温度稳定性 & 好 & 差（天然双折射漂移） \\
\end{tabular}
\tabbot
\normalsize
\end{center}

课件的结论值得原样记住：\textbf{纵向给 KDP/KD$^*$P 做 Q 开关}（$V_\pi$ 由材料定死，
例如 KDP 约 $7.5\ \mathrm{kV}$，但孔径受透明电极限制，简单且温度稳）；
\textbf{横向给 LiNbO$_3$ 做高速调制器}（长晶体 + 窄电极把 $V_\pi$ 压到几百伏以下，
代价是偏置点漂移，需要主动控制）。

%=====================================================================
\hdr{§7}{从相位差到光强：调制器怎么工作}

\subhdr{7.1 装置：为什么偏振器要放 $45^\circ$（课件 p37--p38）}

一条标准的振幅调制链路只有四个块：\textbf{起偏器}（产生干净的线偏振）→ \textbf{电光晶体}
（电压改变两个分量的相位差）→（可选）\textbf{$\lambda/4$ 波片或直流偏置}（把工作点挪到线性区）→
\textbf{检偏器}（把相位变化转成强度变化）。

\begin{center}
\begin{tikzpicture}[font=\footnotesize,x=1cm,y=1cm]
\draw[->,red!65!black,thick] (0,0) -- (12.6,0);
% 起偏器
\draw[fill=bandcream,draw=petal] (1.2,-0.8) rectangle (1.45,0.8);
\draw[<->,violet,thick] (1.05,-0.55) -- (1.6,0.55);
\node[align=center] at (1.33,-1.2) {\scriptsize 起偏器};
% 晶体
\draw[draw=violet,thick,fill=bandviolet] (3.0,-0.9) rectangle (5.8,0.9);
\draw[dashed,inkgray] (3.0,0) -- (5.8,0);
\draw[dashed,inkgray] (3.0,-0.75) -- (5.8,0.75);
\node[align=center,text=violet] at (4.4,2.0) {\scriptsize 晶体主轴：$45^\circ$ 于起偏方向};
\node[align=center] at (4.4,1.25) {\scriptsize 两分量等幅，延迟 $\Gamma(V)$};
\draw[->,violet] (4.4,1.65) -- (4.15,0.75);
\node[align=center] at (4.4,-1.3) {\scriptsize Pockels / Kerr 盒};
% 波片
\draw[fill=bandcream,draw=petal] (6.8,-0.8) rectangle (7.05,0.8);
\node[align=center] at (6.93,-1.2) {\scriptsize $\lambda/4$ 偏置};
% 检偏器
\draw[fill=bandcream,draw=petal] (9.2,-0.8) rectangle (9.45,0.8);
\draw[<->,violet,thick] (9.05,0.55) -- (9.6,-0.55);
\node[align=center] at (9.33,-1.2) {\scriptsize 检偏器（正交）};
\draw[->,red!65!black,thick] (10.6,0) -- (11.6,0) node[below,text=black]{\scriptsize 探测器};
\node[align=center] at (11.0,0.55) {\scriptsize $T=\sin^{2}\dfrac{\Gamma}{2}$};
\end{tikzpicture}
\end{center}

{\small\color{inkgray}图 3：振幅调制链路。晶体主轴与起偏方向成 $45^\circ$ 是关键（§5.2 的 box），
否则一个分量拿不到光、延迟无法转成强度。}

\subhdr{7.2 透射公式的推导}

沿用图 3 的方位：起偏方向取 $x$，晶体主轴在 $\pm45^\circ$，检偏器在 $y$（正交）。

\textbf{第一步，入晶体时等幅。}$E_x=1$ 在晶体坐标（$f$ 快轴与 $s$ 慢轴）上的分解是
$E_f=E_s=1/\sqrt2$，两者同相。

\textbf{第二步，出晶体时有相位差。}经过长度 $L$，慢轴多走 $\Gamma$：
$E_f=1/\sqrt2$，$E_s=\frac{1}{\sqrt2}e^{i\Gamma}$。

\textbf{第三步，投影到检偏器。}$f$ 轴在 $+45^\circ$、$s$ 轴在 $-45^\circ$，它们投到 $y$ 方向分别带
$+1/\sqrt2$ 与 $-1/\sqrt2$ 的系数（符号相反是关键）：
\begin{equation}
E_y=\frac{1}{\sqrt2}\cdot\frac{1}{\sqrt2}
-\frac{1}{\sqrt2}\cdot\frac{1}{\sqrt2}e^{i\Gamma}
=\frac{1}{2}\left(1-e^{i\Gamma}\right).
\end{equation}

\textbf{第四步，取模方。}
\begin{equation}
T=\frac{I_{\text{out}}}{I_{\text{in}}}
=\left|E_y\right|^2=\frac{1}{4}\left|1-e^{i\Gamma}\right|^2
=\frac{1}{4}\left(1+1-2\cos\Gamma\right)
=\frac{1-\cos\Gamma}{2}
=\sin^{2}\frac{\Gamma}{2}.
\end{equation}

于是（课件 p30 的式子）
\begin{equation}
\label{eq:T}
\boxed{\;T=\sin^{2}\frac{\Gamma}{2}=\sin^{2}\!\left(\frac{\pi V}{2V_\pi}\right)\;}
\end{equation}
最后一步用了 $\Gamma=\pi V/V_\pi$，即 (\ref{eq:VpiV})。

\aside{若检偏器与起偏器\emph{平行}而不是正交，同样的推导给出 $T=\cos^2(\Gamma/2)=1-\sin^2(\Gamma/2)$：
两条曲线互补。工程上无所谓哪种更好——但「正交 + 直流偏置 $\pi/2$」是最常见的一种，因为它零电压时全黑，
暗态很干净。}

\subhdr{7.3 三个工作点（课件 p39--p40）}

\begin{center}
\includegraphics[width=0.72\linewidth]{fig-l4/fig2-transmission.png}

{\small\color{inkgray}图 4：$T$--$V$ 曲线（$T=\sin^2\frac{\pi V}{2V_\pi}$）。三个点：
$V=0$ 全黑（null）、$V=V_\pi/2$ 半分（quadrature，也是线性区的中点）、$V=V_\pi$ 全透（peak）。
每过一个 $V_\pi$，曲线重复一次。}
\end{center}

\begin{itemize}
\item \textbf{Null（$V=0$，$T=0$）}：全黑。若在这里直接叠加小交流信号，输出以\emph{两倍}频率振荡
      （因为 $\sin^2$ 在零点附近是偶函数），并产生强烈的谐波失真——不能用来做模拟调制。
\item \textbf{Quadrature / Q 点（$V=V_\pi/2$，$T=0.5$）}：曲线的\emph{拐点}，斜率最大。
      在这里叠加小信号，输出忠实地复制输入波形，失真最小。方法是加 $\lambda/4$ 波片（块状器件）
      或加一个 $V_\pi/2$ 的直流偏置（集成器件）。
\item \textbf{Peak（$V=V_\pi$，$T=1$）}：偏振面正好转过 $90^\circ$，全部透过正交检偏器。
\end{itemize}

注意「null 点两倍频」与 §3.1 里 Kerr 盒「AC 驱动两倍频」是\emph{不同}的两件事：
前者是 $\sin^2$ 在零点附近的偶对称（Pockels 盒也有），后者是 $\Delta n\propto E^2$ 的偶对称（只有 Kerr 有）。
作业里若要区分，就抓住「$\Delta n$ 本身是否反号」这一条。

\subhdr{7.4 线性调制区（课件 p41）}

在 Q 点附近令 $V=V_\pi/2+v(t)$，$|v|\ll V_\pi$。记 $x=\frac{\pi v}{2V_\pi}$，则
\begin{equation}
T=\sin^{2}\!\left(\frac{\pi}{4}+x\right)=\frac{1+\sin 2x}{2}
\approx\frac{1}{2}+x
=\frac{1}{2}+\frac{\pi}{2}\,\frac{v(t)}{V_\pi}.
\end{equation}
（用了 $\sin^{2}(\pi/4+x)=\frac{1+\sin2x}{2}$ 与 $\sin2x\approx2x$。）

\textbf{小例}：$V_\pi=4\ \mathrm{kV}$，加 $v=0.1\ \mathrm{kV}$ 的小信号。
\begin{equation}
\Delta T\approx\frac{\pi}{2}\cdot\frac{0.1}{4}\approx0.039\quad(\text{即 }3.9\%).
\end{equation}
精确值：$T(2.1\ \mathrm{kV})-T(2.0\ \mathrm{kV})=0.0392$，与线性近似差不到 $1\%$。
这说明「小信号线性化」在 $|v|\lesssim0.1V_\pi$ 时已经很好用。

代价也有两条：\textbf{偏置漂移}（温度、电荷俘获会让 Q 点跑掉，必须主动反馈，§7.6）与
\textbf{大信号非线性}（$v$ 一大，$\sin2x$ 的三次项就出来，产生谐波）。

\subhdr{7.5 相位调制器与马赫--曾德尔（课件 p42）}

\textbf{纯相位调制}：撤掉检偏器，只留晶体。此时晶体的一个分量被延迟
\begin{equation}
\Delta\phi(V)=\frac{\pi V}{V_\pi},
\end{equation}
出射仍是线偏振（若偏振沿主轴），只是相位被电压调制。这就是相位调制器（PM）。

\textbf{马赫--曾德尔调制器（MZI）}：把光分成两臂，只在一臂上加相位调制，再合起来：
\begin{equation}
T_{\text{MZI}}=\cos^{2}\!\left(\frac{\Delta\phi}{2}\right)
=\cos^{2}\!\left(\frac{\pi V}{2V_\pi}\right).
\end{equation}
结构上的好处是：两臂的相位噪声共模抵消，而且可以用推挽驱动（一臂 $+\Delta\phi$、另一臂 $-\Delta\phi$），
于是所需电压减半、啁啾（chirp）也小。这是光纤通信里最主流的调制器结构。

\subhdr{7.6 工程难点与「速度极限到底在哪」（课件 p32、p35）}

\begin{center}
\small
\tabtop
\begin{tabular}{@{}p{4.0cm}p{10.0cm}@{}}
\thead{难点} & \thead{现象与对策} \\
\tabmid
温度稳定 & 天然双折射随温度漂移 → 偏置点漂；需温控或补偿 \\
\tabmid
压电振铃 & 电压引起机械应变 → 声学共振；限制带宽，靠安装设计阻尼 \\
\tabmid
光学损伤 & 高功率密度损伤晶体；LiNbO$_3$ 高功率下有光折变；需功率管理 \\
\tabmid
驱动电子学 & 需要高压高速驱动；$>10\ \mathrm{GHz}$ 要用行波电极 \\
\tabmid
迟滞与漂移 & 电荷俘获 → 偏置漂移；需主动偏置控制 \\
\end{tabular}
\tabbot
\normalsize
\end{center}

课件 p35 的第 3 问是这讲最值得记的一条「概念纠偏」：\textbf{Pockels 效应本身是飞秒级的，
限制调制速度的从来不是材料}。真正卡住带宽的是五件事：

（1）\textbf{RC 时间常数}：块状电极与晶体构成电容（几 pF），配 $50\ \Omega$ 驱动就是几 GHz 的低通；
（2）\textbf{速度失配}：$>10\ \mathrm{GHz}$ 用行波电极时，微波折射率（LiNbO$_3$ 约 $4$--$5$）与光折射率（约 $2.2$）不同，
调制波「追不上」光波，作用长度受限；
（3）\textbf{微波损耗}：电极电阻与介质损耗随频率上升；
（4）\textbf{压电振铃}：逆压电效应激发声学共振，$100\ \mathrm{MHz}$--$1\ \mathrm{GHz}$ 就有影响；
（5）\textbf{驱动源}：$>40\ \mathrm{GHz}$、又要几十伏摆幅，本身就难做。

现状：商用 LiNbO$_3$ 调制器 $40$--$100\ \mathrm{GHz}$；薄膜铌酸锂与等离子体结构在实验室已到 $100\ \mathrm{GHz}$ 以上。

\subhdr{7.7 Q 开关：把电光效应当「腔门」用（课件 p36 的应用目录）}

课件把 Q 开关列在应用里但没展开讲原理，这里补到能读懂器件的程度。

\textbf{为什么要 Q 开关。}腔的品质因数 $Q$ 大致是「腔里存的能量 ÷ 每周期损耗的能量」。$Q$ 高，振荡阈值低，
一旦起振反转粒子数就被迅速消耗，输出是连续的、功率受限于阈值；$Q$ 低，阈值高，泵浦可以先把
反转粒子数\emph{攒}到远高于阈值，然后突然降低损耗，一次把储能放出来——这就是巨脉冲（脉宽 ns 级，
峰值功率可到 GW）。

\textbf{电光 Q 开关怎么接。}腔内放「起偏器 + Pockels 盒」，盒的主轴与起偏方向成 $45^\circ$。

\begin{enumerate}
\item \textbf{关闭态（加电压）}：盒调到\emph{单程} $\lambda/4$（即 $V=V_\pi/2$）。光往一趟走，两个分量差 $90^\circ$；
      经端镜反射再走一趟，延迟\emph{继续累积}到 $180^\circ$——等效于一片半波片，偏振面转过 $90^\circ$，
      回到起偏器时被挡住。腔损耗高，$Q$ 低，不振荡。
\item \textbf{泵浦}：反转粒子数在这一段时间里堆积到很高。
\item \textbf{打开（撤电压，ns 级）}：延迟消失，损耗骤降，$Q$ 变高，巨脉冲在几十个往返内建立。
\end{enumerate}

\textbf{两个值得注意的细节。}

（i）为什么是 $V_\pi/2$ 而不是 $V_\pi$：因为光要\emph{走两趟}。这与 §6.2 里「单程」的 $V_\pi$ 定义不同，
买器件时会看到「四分之一波电压」这个规格，就是为这个用法标的。

（ii）为什么往返延迟会\emph{累积}：晶体的主轴由\emph{电场}确定，不随光的传播方向反号，
所以回程看到的还是同一对快慢轴、同一个 $\Delta n$，相位差自然相加。
\textbf{记住这句话}——§12 会看到 Faraday 旋转的「累积」有完全相同的原因（场定的方向不随 $k$ 反号），
而自然旋光的「抵消」正因为它依赖的是介质的\emph{手性}而不是外场。

\textbf{与 Kerr 盒 Q 开关的对比}：Kerr 盒也能做（大孔径、不挑极性），但所需电压高、液体维护麻烦，
现代固体激光器几乎都用 Pockels 盒（KD$^*$P、BBO 或 LiNbO$_3$）。

%=====================================================================
\hdr{§8}{课件里的六道例题（教学示范）}

{\small\color{inkgray}这六题课件已给完整解答，所以按\emph{教学示范}处理：正文给出完整演算，
不作为需要先做的独立练习。作业三题（§16）用的是另一组数字，不给步骤。}

\subhdr{8.1 例题 1：KDP 横向盒的半波电压（课件 p43--p45）}

\textbf{已知}：KDP 横向盒，$L=2\ \mathrm{cm}$，电极间距 $d=5\ \mathrm{mm}$，$r_{63}=10.5\times10^{-12}\ \mathrm{m/V}$，
$n_o=1.51$，$\lambda=1550\ \mathrm{nm}$。求 $V_\pi$。

\textbf{解}：横向几何用 $\Gamma=\frac{2\pi}{\lambda}n_o^3r_{63}\frac{L}{d}V$，令 $\Gamma=\pi$：
\begin{equation}
V_\pi=\frac{\lambda}{2n_o^3r_{63}}\cdot\frac{d}{L}
=\frac{1.55\times10^{-6}}{2\times1.51^3\times10.5\times10^{-12}}\cdot\frac{0.005}{0.02}.
\end{equation}
先算材料因子：$2\times1.51^3\times10.5\times10^{-12}\approx7.23\times10^{-11}\ \mathrm{m/V}$，于是
\begin{equation}
\frac{\lambda}{2n_o^3r_{63}}\approx\frac{1.55\times10^{-6}}{7.23\times10^{-11}}\approx2.14\times10^{4}\ \mathrm{V};
\qquad
V_\pi\approx2.14\times10^{4}\times0.25\approx5.36\times10^{3}\ \mathrm{V}.
\end{equation}

\textbf{答}：$V_\pi\approx5.4\ \mathrm{kV}$。

\textbf{两点补充。}（i）把 $\lambda$ 从 $1550\ \mathrm{nm}$ 换到 $633\ \mathrm{nm}$，$V_\pi$ 按正比缩小到约 $2.2\ \mathrm{kV}$——
半波电压与波长成正比。（ii）这里的 $n_o^3$ 与 $r_{63}$ 都取自 KDP 在近红外的值；
作业第 1 题给的是 $n=1.5$、$r=10.6\ \mathrm{pm/V}$ 的简化值，道理一样。

\subhdr{8.2 例题 2：液体 Kerr 盒的四分之一波电压（课件 p46--p48）}

\textbf{已知}：$K=2.2\times10^{-14}\ \mathrm{m/V^2}$，$d=1\ \mathrm{cm}$，$L=5\ \mathrm{cm}$，
$\lambda=632.8\ \mathrm{nm}$，要求 $\Gamma=\pi/2$。求 $V$。

\textbf{解}：由 (\ref{eq:gamma-kerr})，
\begin{equation}
\Gamma=2\pi KL\left(\frac{V}{d}\right)^2
\;\Longrightarrow\;
\frac{V}{d}=\sqrt{\frac{\Gamma}{2\pi KL}} .
\end{equation}
分母：$2\pi KL=2\pi\times2.2\times10^{-14}\times0.05=6.91\times10^{-15}$。于是
\begin{equation}
\frac{V}{d}=\sqrt{\frac{\pi/2}{6.91\times10^{-15}}}\approx\sqrt{2.27\times10^{14}}\approx1.51\times10^{7}\ \mathrm{V/m},
\qquad
V\approx1.51\times10^{7}\times0.01\approx1.51\times10^{5}\ \mathrm{V}.
\end{equation}

\textbf{答}：$V\approx1.5\times10^{5}\ \mathrm{V}=150\ \mathrm{kV}$。

\textbf{为什么要出这道题}：它让你对「Kerr 效应弱」有数字感。同一个 $L$、$d$ 下，硝基苯
（$K$ 大约 $100$ 倍）只要约 $1.5\ \mathrm{kV}$ 量级的场强就能做到，而这里的 $K$ 比硝基苯小两个数量级，
电压就要高一个数量级到 $150\ \mathrm{kV}$。顺带注意：延迟与 $\lambda$ 无关（$\lambda$ 在 (\ref{eq:gamma-kerr}) 里消掉了），
题目给出 $632.8\ \mathrm{nm}$ 只是为提示光源。

\subhdr{8.3 例题 3：Pockels 盒的强度调制（课件 p49--p52）}

\textbf{已知}：Pockels 盒在正交偏振器之间，$V_\pi=4\ \mathrm{kV}$，输入 $I_{\text{in}}=2\ \mathrm{mW}$。
求 (a) $I_{\text{out}}(V)$；(b) 透过最大值 $50\%$ 的电压；(c) 该电压下的输出功率。

\textbf{解 (a)}：把 $\Gamma=\pi V/V_\pi$ 代入 (\ref{eq:T})：
\begin{equation}
I_{\text{out}}(V)=I_{\text{in}}\sin^{2}\!\left(\frac{\pi V}{2V_\pi}\right).
\end{equation}

\textbf{解 (b)}：令 $I_{\text{out}}=0.5I_{\text{in}}$（注意 $I_{\max}=I_{\text{in}}$，因为 $\sin^2$ 最大值为 $1$）：
\begin{equation}
\sin^{2}\!\left(\frac{\pi V}{2V_\pi}\right)=0.5
\;\Longrightarrow\;
\frac{\pi V}{2V_\pi}=\frac{\pi}{4}
\;\Longrightarrow\;
V=\frac{V_\pi}{2}=2\ \mathrm{kV}.
\end{equation}

\textbf{解 (c)}：$I_{\text{out}}=0.5\times2\ \mathrm{mW}=1\ \mathrm{mW}$。

\textbf{答}：(a) $I_{\text{in}}\sin^{2}(\pi V/2V_\pi)$；(b) $2\ \mathrm{kV}$；(c) $1\ \mathrm{mW}$。

\textbf{一句话点评}：这题就是 §7.3 的 Q 点。$50\%$ 处既是「一半光」，也是线性调制的最佳偏置点。

\subhdr{8.4 例题 4--6：磁光三题（课件 p79--p81）}

\textbf{例题 4（TGG 转角，p79）}：$L=5\ \mathrm{cm}$，$B=0.8\ \mathrm{T}$，$\vdt=40\ \mathrm{rad/(T\cdot m)}$。
\begin{equation}
\theta=\vdt BL=40\times0.8\times0.05=1.6\ \mathrm{rad}=91.7^\circ .
\end{equation}

\textbf{例题 5（隔离器长度，p80）}：要 $45^\circ$，$\vdt=25\ \mathrm{rad/(T\cdot m)}$，$B=0.5\ \mathrm{T}$。
\begin{equation}
L=\frac{\theta}{\vdt B}=\frac{\pi/4}{25\times0.5}=\frac{0.7854}{12.5}\approx0.0628\ \mathrm{m}\approx6.3\ \mathrm{cm}.
\end{equation}

\textbf{例题 6（两材料比较，p81）}：$B=0.3\ \mathrm{T}$，要 $45^\circ$；$\vdt_A=30$，$\vdt_B=120\ \mathrm{rad/(T\cdot m)}$。
\begin{equation}
L_A=\frac{0.7854}{30\times0.3}\approx0.087\ \mathrm{m}=8.7\ \mathrm{cm},
\qquad
L_B=\frac{0.7854}{120\times0.3}\approx0.022\ \mathrm{m}=2.2\ \mathrm{cm}.
\end{equation}
$\vdt$ 大 $4$ 倍，长度就短 $4$ 倍：\textbf{大 Verdet 常数的材料（如 TGG、Ce:YIG）是紧凑隔离器的前提}（§15.4）。

%=====================================================================
\hdr{§9}{自然旋光：一个必须拿来对照的「互易」样本}

{\small\color{inkgray}课件 p2 的课程计划把「自然旋光（Fresnel 圆双折射、糖度计）」列为第 3 部分，
但 p53 直接从电光跳到磁光，正课页里没有展开。它在 L2（旋光）与 L3（糖度计）已经讲过，
这里只补到「足以和 Faraday 对照」的程度——因为 §12 的非互易性离开这个对照就讲不清。}

\subhdr{9.1 现象与 Biot 定律}

把一片偏振器与一片检偏器正交放置，中间插一根装糖溶液的管子，视场会亮起来；转动检偏器到某个角度
又重新变黑——说明溶液把偏振面\emph{转了}一个角。这个角与浓度、与光程都成正比：
\begin{equation}
\label{eq:biot}
\theta=[\alpha]\,c\,L ,
\end{equation}
$[\alpha]$ 是比旋光度（由材料、波长、温度决定），$c$ 是浓度，$L$ 是光程。这就是糖度计（saccharimetry）
的全部原理：\textbf{测一个角度，反推一个浓度}。

\subhdr{9.2 Fresnel 的圆双折射模型（为什么会转）}

Fresnel 的解释只有一句话：\textbf{手性介质对右旋与左旋圆偏振的折射率不同}（圆双折射），
而线偏振是这两者的等幅叠加，两者走完同一段长度后相位不同，合出来的线偏振自然转了个方位。

\textbf{推导。}记 $\hat e_\pm=(\hat x\pm i\hat y)/\sqrt2$，两个模的折射率为 $n_\pm$。入射时等幅同相：
\begin{equation}
\vec E(0)=\frac{E_0}{\sqrt2}\left(\hat e_+ +\hat e_-\right).
\end{equation}
走完 $L$ 之后各自累积 $2\pi n_\pm L/\lambda$：
\begin{equation}
\vec E(L)=\frac{E_0}{\sqrt2}\left(\hat e_+\,e^{i2\pi n_+L/\lambda}+\hat e_-\,e^{i2\pi n_-L/\lambda}\right)
=e^{i\bar\Phi}\,\frac{E_0}{\sqrt2}\left(\hat e_+\,e^{i\delta/2}+\hat e_-\,e^{-i\delta/2}\right),
\end{equation}
其中 $\bar\Phi=\pi(n_++n_-)L/\lambda$，$\delta=2\pi(n_+-n_-)L/\lambda$。把 $\hat e_\pm$ 展开：
\begin{equation}
\hat e_+\,e^{i\delta/2}+\hat e_-\,e^{-i\delta/2}
=\sqrt2\left[\hat x\cos\frac{\delta}{2}-\hat y\sin\frac{\delta}{2}\right].
\end{equation}
右边是一个\emph{线}偏振，方位从 $\hat x$ 转了 $-\delta/2$。所以
\begin{equation}
\label{eq:oa}
\boxed{\;\theta=\frac{\pi}{\lambda}\left(n_- -n_+\right)L=\frac{\pi}{\lambda}\,\Delta n_c\,L\;}
\qquad \Delta n_c\equiv n_- -n_+ .
\end{equation}

形式与 (\ref{eq:gamma}) 一样：\textbf{相位差除以 2 就是方位角的转角}。这不是巧合——
把 (\ref{eq:gamma}) 与 (\ref{eq:oa}) 并排看，$\Gamma\leftrightarrow\delta$、$2\theta\leftrightarrow\delta$，
「转一个角 $\theta$」与「两个圆分量差 $2\theta$」是同一件事的两种说法。

\subhdr{9.3 为什么它是互易的（\S12 的伏笔）}

关键在 $\Delta n_c$ 的来源：它来自介质的\emph{螺旋结构}，而「右旋/左旋」这个标签是相对\emph{传播方向}定的。
光反射回来时传播方向反了，于是同一束物理光在新的传播方向上对应的圆手性\emph{交换}：
去程「快」的那个模，回程变「慢」。相位差被原路退回，净转角为零。这就是\textbf{互易}。

微观上说得更清楚：自然旋光源于介质的空间色散（效应正比于波矢 $\vec k$，$\vec k$ 反号则效应反号），
而 Faraday 效应源于磁化（效应正比于 $\vec B$，与 $\vec k$ 无关）。这句话是全讲最重要的一句对照，§12 会把它算实。

%=====================================================================
\hdr{§10}{磁光效应：一张表看清四类}

\subhdr{10.1 定义与两个基本量（课件 p53--p54）}

磁光（MO）效应是「光与\emph{磁化}介质相互作用，使偏振态、强度或相位改变」的一族现象。
现象学上它们都表现为两个量：

\begin{itemize}
\item \textbf{磁圆双折射（MCB）}：RCP 与 LCP 的\emph{折射率}不同 → 偏振面\emph{转}（Faraday，§11）；
\item \textbf{磁圆二色性（MCD）}：RCP 与 LCP 的\emph{吸收}不同 → 偏振态变\emph{椭圆}（§14.2）。
\end{itemize}

课件 p54 有一句话必须原样记住：\textbf{磁光效应的定义特征是「非互易」——一次传播中累积的旋转，
不会因为反射或反向传播而抵消}。这句话就是 §12 的全部内容，也是隔离器能工作的唯一理由。

\subhdr{10.2 介电张量：非对角元就是磁光效应（课件 p57）}

各向同性介质的介电张量是对角的。磁化之后出现非对角元：
\begin{equation}
\underline{\underline{\varepsilon}}=\varepsilon_0
\begin{pmatrix}
\varepsilon & \varepsilon_{xy} & \varepsilon_{xz}\\
-\varepsilon_{xy} & \varepsilon & -\varepsilon_{yz}\\
-\varepsilon_{xz} & \varepsilon_{yz} & \varepsilon
\end{pmatrix}
=\varepsilon_0
\begin{pmatrix}
\varepsilon & iQ_z & iQ_y\\
-iQ_z & \varepsilon & -iQ_x\\
-iQ_y & iQ_x & \varepsilon
\end{pmatrix},
\end{equation}
其中 $Q_j$ 是磁光常数（与磁化强度或外场成正比）。\textbf{对角元}给出无场时的各向同性响应，
\textbf{非对角元}给出全部磁光效应；这就是为什么「磁光效应弱」等价于「$\varepsilon_{xy}$ 小」
（$Q\sim10^{-5}$--$10^{-3}$，见 §15.2）。

\textbf{三种几何给出三种效应}（课件 p57 的分类）：磁化平行于光的传播方向 → 透射时是 Faraday 效应、
反射时是极向 MOKE；磁化垂直于传播方向 → 横向（Voigt/Cotton--Mouton）效应。
课件把第三类写成「Absorption → Voigt Effect」，这个标签有问题：Voigt 效应是\emph{双折射}（相位），
属于吸收的是 MCD。核对见附录 B 第 5 条。

\subhdr{10.3 名词表（课件 p56）}

\begin{center}
\small
\tabtop
\begin{tabular}{@{}p{3.6cm}p{7.4cm}p{3.0cm}@{}}
\thead{名称} & \thead{内容} & \thead{几何} \\
\tabmid
Faraday 旋转 & 透射光的偏振面旋转 & $\vec B\parallel\vec k$ \\
\tabmid
MO Kerr 效应（MOKE） & 反射光的偏振态改变（旋转 + 椭圆率） & 反射，表面敏感 \\
\tabmid
Voigt 效应 & 磁线双折射（气体） & $\vec B\perp\vec k$ \\
\tabmid
Cotton--Mouton 效应 & 磁线双折射（液体；与 Voigt 同一现象） & $\vec B\perp\vec k$ \\
\tabmid
MCD & LCP 与 RCP 的吸收差 & $\vec B\parallel\vec k$ \\
\tabmid
Zeeman 效应 & 磁场使谱线分裂 & 原子能级 \\
\end{tabular}
\tabbot
\normalsize
\end{center}

\subhdr{10.4 年表（课件 p55）}

\begin{center}
\small
\tabtop
\begin{tabular}{@{}p{2.0cm}p{6.0cm}p{6.0cm}@{}}
\thead{年份} & \thead{发现} & \thead{人物} \\
\tabmid
1845 & Faraday 效应 & Michael Faraday \\
\tabmid
1877 & （磁光）Kerr 效应 & John Kerr \\
\tabmid
1902 & Voigt 效应 & Woldemar Voigt \\
\tabmid
1907 & Cotton--Mouton 效应 & Aimé Cotton \& Henri Mouton \\
\tabmid
1957 & 首次磁光记录（MnBi） & — \\
\tabmid
1980s & 磁光盘商品化 & — \\
\end{tabular}
\tabbot
\normalsize
\end{center}

注意这里有两个 Kerr：\textbf{1875 年的电光 Kerr 效应}（§3，二次电光）与
\textbf{1877 年的磁光 Kerr 效应}（§13，反射）。同名不同事，考试里常被拿来设陷阱。

\subhdr{10.5 磁性类型决定磁光响应（课件 p68）}

\begin{center}
\small
\tabtop
\begin{tabular}{@{}p{3.2cm}p{3.4cm}p{3.6cm}p{3.4cm}@{}}
\thead{类型} & \thead{磁化率 $\chi$} & \thead{$M$--$H$ 关系} & \thead{非对角元正比于} \\
\tabmid
顺磁 & 正、小 & 线性 & 外场 $H$ \\
\tabmid
抗磁 & 负、小 & 线性 & 外场 $H$ \\
\tabmid
铁磁 & 正、大 & 非线性（低场饱和） & 磁化强度 $M$ \\
\end{tabular}
\tabbot
\normalsize
\end{center}

所有材料都有磁光活性，强度差别在于磁性：铁磁材料的磁光常数比顺磁/抗磁高几个数量级，
因为它的非对角元正比于\emph{自发磁化}而不是外加场的微小极化。这条是 §15.1 选材料的依据。

%=====================================================================
\hdr{§11}{Faraday 旋转：$\theta=\mathcal{V}BL$}

\subhdr{11.1 公式与 Verdet 常数（课件 p58）}

\begin{equation}
\label{eq:faraday}
\boxed{\;\theta=\vdt BL\;}
\end{equation}
$\vdt$ 是 Verdet 常数（单位 $\mathrm{rad/(T\cdot m)}$，符号与旋向约定有关），
$B$ 是磁感应强度沿光传播方向的分量，$L$ 是光程。更一般地
\begin{equation}
\theta=\vdt\int \vec B\cdot\mathrm d\vec l ,
\end{equation}
这一行解释了 §12 里的一切：\textbf{转角由 $\vec B$ 与路径的相对取向决定}，所以光反向走时
$\mathrm d\vec l$ 反号而 $\vec B$ 不反号，积分值按同样大小累积（不是抵消）。

Faraday 1845 年发表。它在抗磁、顺磁、铁磁材料里都有；铁磁的钇铁石榴石（YIG）
在可见与近红外的 Verdet 常数极大，是器件材料（§15.4）。

\subhdr{11.2 从介电张量看：旋转与椭圆率是两个分量（课件 p59）}

只考虑传播方向（取 $z$）上的两个圆模。把 §10.2 的张量作用到 $\vec E_\pm=E_x\pm iE_y$ 上：
\begin{equation}
D_\pm=\varepsilon_0\left(\varepsilon \mp i\varepsilon_{xy}\right)E_\pm
\;\Longrightarrow\;
n_\pm^2=\varepsilon \mp i\varepsilon_{xy} .
\end{equation}
于是
\begin{equation}
n_\pm\approx n\left(1\mp\frac{i\varepsilon_{xy}}{2\varepsilon}\right),\qquad n=\sqrt\varepsilon ,
\end{equation}
两个圆模的折射率之差 $n_+-n_-=-i\varepsilon_{xy}/n$ \textbf{是复数}，两个部分各管一件事：

\begin{itemize}
\item 虚部（等价于 $\varepsilon_{xy}$ 的实部）→ 相位差 → \textbf{旋转 $\theta$}（圆双折射，MCB）；
\item 实部（等价于 $\varepsilon_{xy}$ 的虚部）→ 振幅差 → \textbf{椭圆率 $\psi$}（圆二色性，MCD）。
\end{itemize}

课件 p59 把它写成
\begin{equation}
\label{eq:eps-rot}
\theta=-\frac{\omega}{2c}\,\mathrm{Im}\!\left(\frac{\varepsilon_{xy}}{\varepsilon}\right),
\qquad
\psi=\frac{\omega}{2c}\,\mathrm{Re}\!\left(\frac{\varepsilon_{xy}}{\varepsilon}\right),
\end{equation}
它给出的是\textbf{单位长度}的转角与椭偏率（再乘 $L$ 才是总角），系数里还含一个折射率因子，
取决于 $\varepsilon$ 取绝对还是相对值。结构（Im 管旋转、Re 管椭圆）是对的，量级也是对的；
但\textbf{答题与做题一律回到 (\ref{eq:faraday}) 这个可测量的形式}。核对见附录 B 第 1 条。

\aside{为什么「旋转」与「椭圆率」共用同一个 $\varepsilon_{xy}$ 的两个部分？因为它们是同一个复数折射率差的实部与虚部——
这与「复折射率 $n+i\kappa$ 的实部管相位、虚部管吸收」是同一件事。Kramers--Kronig 关系进一步说明：
只要一个频率上有 CD，就必然在附近有对应的 CB（课件 p73 的表）。}

\subhdr{11.3 Verdet 常数的实际数值}

\begin{center}
\small
\tabtop
\begin{tabular}{@{}p{6.0cm}p{4.4cm}p{5.0cm}@{}}
\thead{材料} & \thead{$\vdt$，$\mathrm{rad/(T\cdot m)}$} & \thead{备注} \\
\tabmid
TGG（铽镓石榴石） & $\sim40$（$633\ \mathrm{nm}$） & 可见--近红外主力 \\
\tabmid
普通玻璃 & $\sim10$--$20$ & 便宜、可做大孔径 \\
\tabmid
YIG（钇铁石榴石） & 近红外极大 & 铁磁，磁光器件 \\
\tabmid
Ce:YIG & 比 YIG 高一个量级 & 集成隔离器 \\
\end{tabular}
\tabbot
\normalsize
\end{center}

量级感：$\vdt\sim20\ \mathrm{rad/(T\cdot m)}$、$B=1\ \mathrm{T}$ 时，要转 $45^\circ$ 需要
$L=\theta/(\vdt B)=0.785/20\approx4\ \mathrm{cm}$。这就是为什么隔离器里的磁光棒总有几厘米长，
而高 Verdet 材料（Ce:YIG）能把器件做到毫米级（§15.4）。

%=====================================================================
\hdr{§12}{非互易性：为什么往返是 $2\theta$}

\subhdr{12.1 问题的提法}

把一束偏振光送进介质，再从端面反射回来，重新穿过同一段介质。问出射的偏振面转了多少度。

两个候选答案，差一个符号：
\begin{itemize}
\item \textbf{互易（光学活性）}：去程 $+\theta$，回程 $-\theta$，合计 $0$——旋光「退回去」。
\item \textbf{非互易（Faraday）}：去程 $+\theta$，回程\emph{还是} $+\theta$，合计 $2\theta$。
\end{itemize}

课件 p76 的 FAQ 给的是后者，而且理由说得很准：\textbf{旋转方向由磁场方向定，不由光的传播方向定}。
这一节把这个理由算实。

\subhdr{12.2 严格论证：张量不变，相位照加}

全部论证只用一件事：\textbf{介质的介电张量由 $\vec B$ 决定，而 $\vec B$ 在光反射时并不反号}。

\textbf{第一步（正程）。}取固定的实验室坐标 $(\hat x,\hat y)$。在 $z=0$ 处入射线偏振，
分解成两个圆模，等幅同相。走到 $z=L$（端镜处）：
\begin{equation}
\vec E(L)=\frac{E_0}{\sqrt2}\left(\hat e_+\,e^{ik_+L}+\hat e_-\,e^{ik_-L}\right),
\qquad k_\pm=\frac{\omega}{c}n_\pm .
\end{equation}
由 §9.2 的算法，偏振面在实验室坐标里转了 $\theta=\frac{\pi}{\lambda}(n_- - n_+)L$。

\textbf{第二步（反射）。}理想端镜把两个分量都乘上 $-1$（切向场反向），\emph{相对}相位不变：
$\vec E'(L)=-\vec E(L)$。它的矢量分量与入射时相同，但传播方向变成 $-\hat z$。

\textbf{第三步（回程）。}反向传播的模式因子是 $e^{-ik_\pm z}$（$k_\pm>0$）。从 $z=L$ 走到 $z=0$，
相位变化是 $-k_\pm(0-L)=+k_\pm L$——与正程的 $+k_\pm L$ \emph{同号同大小}。于是
\begin{equation}
\vec E(0)\propto\frac{E_0}{\sqrt2}\left(\hat e_+\,e^{2ik_+L}+\hat e_-\,e^{2ik_-L}\right),
\end{equation}
两个圆分量之间的相对相位差是正程的两倍：
\begin{equation}
2(k_+-k_-)L \;\Longrightarrow\; \theta_{\text{往返}}=2\theta=2\vdt BL .
\end{equation}


\textbf{一句话总结}：正程与回程的差别只在于「往哪个方向走」，而 $n_\pm$ 由 $\vec B$ 定、
跟往哪个方向走无关。这就是「非互易」的精确含义。

\aside{用 Jones 矩阵写更短：在固定实验室基 $(\hat x,\hat y)$ 下，正程介质的作用是旋转矩阵 $R(\theta)$，
回程介质的作用\emph{还是} $R(\theta)$（不是 $R(-\theta)$），所以往返是 $R(\theta)R(\theta)=R(2\theta)$。
若写成 $R(\theta)R(-\theta)=I$，就是把一个非互易元件误当成了互易元件——这是本节唯一的坑。}

\subhdr{12.3 与自然旋光的对照：为什么那边是 $0$}

同样三步，把介质换成糖溶液或石英（§9）。差别出现在\textbf{$n_\pm$ 是否交换}：

\begin{itemize}
\item \textbf{自然旋光}：$n_\pm$ 的定义依赖「相对传播方向的手性」。回程时传播方向反号，
      同一个物理螺旋对反向波呈现相反的手性，于是去程快的模在回程变慢：
      $n_\pm^{\text{回}}=n_\mp^{\text{正}}$，相对相位差反号，$\theta_{\text{回}}=-\theta$，合计 $0$。
\item \textbf{Faraday}：$n_\pm$ 的定义依赖 $\vec B$（轴矢量，由线圈电流定）。$\vec B$ 不随 $\vec k$ 反号，
      所以 $n_\pm^{\text{回}}=n_\pm^{\text{正}}$，相对相位差同号，$\theta_{\text{回}}=+\theta$，合计 $2\theta$。
\end{itemize}

微观根源一句话：\textbf{自然旋光来自空间色散（效应 $\propto\vec k$），Faraday 来自磁化（效应 $\propto\vec B$）}。

\begin{center}
\begin{tikzpicture}[font=\footnotesize,x=1cm,y=1cm]
% ---------- (a) Faraday ----------
\node at (3.6,2.0) {\textbf{(a) Faraday：不抵消，$2\theta$}};
\draw[->,inkgray] (0.2,0.9) -- (7.0,0.9);
\draw[draw=violet,thick,fill=bandviolet] (2.4,0.35) rectangle (4.4,1.45);
\node[align=center,text=violet] at (3.4,0.02) {\scriptsize 磁光介质（$\vec B$ 固定）};
\draw[draw=violet,thick] (6.9,0.2) rectangle (7.15,1.6);
\node[align=center] at (7.02,-0.55) {\scriptsize 端镜};
\draw[<->,violet,thick] (0.6,0.65) -- (0.6,1.15); \node at (0.6,1.5) {\scriptsize $0^\circ$};
\draw[<->,violet,thick] (5.0,0.52) -- (5.75,1.28); \node at (5.6,1.6) {\scriptsize $45^\circ$};
\draw[<->,violet,thick] (6.0,0.9) -- (6.5,0.9); \node at (6.0,0.5) {\scriptsize $90^\circ$};
\draw[->,red!65!black] (7.35,0.9) -- (7.9,0.9);
\node[align=center] at (7.62,1.35) {\scriptsize 被挡住}; \draw[red!65!black,thick] (7.5,0.62) -- (7.75,1.18);
\node[align=left,text=inkgray] at (3.9,-1.15)
  {\scriptsize 回程再转 $45^\circ$ → 与起偏方向正交 → \textbf{隔离器原理}};

% ---------- (b) 自然旋光 ----------
\node at (3.6,-2.5) {\textbf{(b) 自然旋光：原路退回，$0$}};
\draw[->,inkgray] (0.2,-3.6) -- (7.0,-3.6);
\draw[draw=petal,thick,fill=bandcream] (2.4,-4.15) rectangle (4.4,-3.05);
\node[align=center] at (3.4,-4.55) {\scriptsize 手性介质（无外场）};
\draw[draw=petal,thick] (6.9,-4.3) rectangle (7.15,-2.9);
\node[align=center] at (7.02,-4.9) {\scriptsize 端镜};
\draw[<->,violet,thick] (0.6,-3.35) -- (0.6,-3.85); \node at (0.6,-3.05) {\scriptsize $0^\circ$};
\draw[<->,violet,thick] (5.0,-3.98) -- (5.75,-3.22); \node at (5.6,-2.95) {\scriptsize $45^\circ$};
\draw[<->,violet,thick] (6.0,-3.6) -- (6.5,-3.6); \node at (6.1,-4.02) {\scriptsize $0^\circ$};
\draw[->,red!65!black] (7.35,-3.6) -- (7.9,-3.6);
\node[align=center] at (7.62,-3.2) {\scriptsize 通过};
\node[align=left,text=inkgray] at (3.9,-5.55)
  {\scriptsize 回程反向抵消 $45^\circ$ → 回到 $0^\circ$ → \textbf{不能做隔离器}};
\end{tikzpicture}
\end{center}

{\small\color{inkgray}图 5：同一段介质，光走两趟。(a) Faraday：转角按 $+45^\circ\to+90^\circ$ 累积；
(b) 自然旋光：$+45^\circ\to0^\circ$ 退回。判据不是「有没有磁」，而是「决定旋向的量随不随 $\vec k$ 反号」。}

\subhdr{12.4 光隔离器的构造（课件 p74）}

图 5(a) 就是光隔离器的全部。三个元件沿光轴排好：
\begin{enumerate}
\item \textbf{起偏器}（方位 $0^\circ$）；
\item \textbf{Faraday 旋转器}（转 $45^\circ$，用高 Verdet 材料 + 永磁体提供 $\vec B$）；
\item \textbf{检偏器}（方位 $45^\circ$，与起偏器成 $45^\circ$）。
\end{enumerate}
\textbf{正向光}：$0^\circ\to45^\circ$，与检偏器一致，通过。
\textbf{返回光}（由下游反射回来）：过检偏器后是 $45^\circ$，再过旋转器再加 $45^\circ$ → $90^\circ$，
与起偏器（$0^\circ$）正交 → 被挡住。

\textbf{关键点}：如果旋转器是「互易」的（例如一片 $\lambda/2$ 波片或一根糖溶液），返回光会退回 $0^\circ$
并顺利穿过，隔离失效。\textbf{所以隔离器必须用非互易元件}——这就是整个 §12 的工程意义。

两个配套名词：\textbf{隔离度}（isolation，返回光被压掉多少 dB，商用 $30$--$60\ \mathrm{dB}$）与
\textbf{插入损耗}（forward insertion loss，通常 $0.5$--$1\ \mathrm{dB}$）。
\textbf{环形器}是它的多端口推广：$1\to2\to3\to1$ 单向循环，靠同一套「非互易旋转 + 偏振分束」。

\subhdr{12.5 关于作业第 2 题 (b) 的措辞}

作业原文是「若光两次穿过棒（往返），总转角是多少？解释为什么结果不是简单地翻倍」。
物理结论明确：\textbf{总转角是 $2\theta$}（课件 p76 的 FAQ 也是这么写的）。所以这句话不能被读成
「答案不是 $2\theta$」。它要你说明的是\textbf{「为什么不能按『反向传播所以反向抵消』的直觉去算」}——
也就是本节 12.2 的三步论证。若作业想考的是「为什么不是 $0$」，那答案就是非互易性本身。

写法上最稳的一句：\emph{The rotation doubles in the laboratory frame because the sense of rotation is
fixed by $\vec B$, not by the direction of propagation; unlike natural optical activity, the return pass
adds another $+\theta$ instead of undoing the first one.} 附录 C 与附录 D 都按这个口径给答案。

% ==== PART7 ====

%=====================================================================
\hdr{§13}{MOKE：反射时的偏振改变}

\subhdr{13.1 定义与表面敏感性（课件 p60）}

磁光 Kerr 效应（MOKE，1877）说的是：\textbf{线偏振光从磁化表面反射后，偏振态改变}（既有旋转 $\theta$，
也有椭圆率 $\psi$）。它与 Faraday 效应同源（都来自 $\varepsilon_{xy}$），区别只在几何：一个透射、一个反射。

\textbf{反射意味着表面敏感}：探测深度约为光的穿透深度，金属多层膜里约 $20\ \mathrm{nm}$。
这是它在薄膜磁学里不可替代的原因——它只"看"表面那二十纳米。

\begin{WarnBox}
课件 p60 有一句话要加限定：「效应正比于\emph{沿光传播方向}的磁化分量」。
这只对\textbf{极向（polar）}几何（磁化垂直于膜面）成立；纵向与横向几何的正比关系分别是
「面内沿入射面的分量」与「面内垂直入射面的分量」（§13.2）。核对见附录 B 第 4 条。
\end{WarnBox}

\subhdr{13.2 三种几何（课件 p61）}

\begin{center}
\begin{tikzpicture}[font=\footnotesize,x=1cm,y=1cm]
% 公共画法：膜面 y=0.55，光从左上到 (x0+1.3,0.55) 再反射到右上。
% 磁化标记画在光束右侧，只标单字母 M；含义见图注与下方表，避免文字压住光路。
% --- (a) polar ---
\node at (1.9,2.25) {\textbf{(a) 极向 polar}};
\draw[thick,violet] (0.0,0.55) -- (3.9,0.55);
\draw[->,red!65!black] (0.1,1.75) -- (1.3,0.56);
\draw[->,red!65!black] (1.3,0.56) -- (2.5,1.75);
\node at (0.35,1.2) {\scriptsize 入};
\node at (2.28,1.2) {\scriptsize 反};
\draw[->,violet,thick] (3.15,0.62) -- (3.15,1.40);
\node at (3.48,1.15) {\scriptsize $\vec M$};

% --- (b) longitudinal ---
\node at (6.9,2.25) {\textbf{(b) 纵向 longitudinal}};
\draw[thick,violet] (5.0,0.55) -- (8.9,0.55);
\draw[->,red!65!black] (5.1,1.75) -- (6.3,0.56);
\draw[->,red!65!black] (6.3,0.56) -- (7.5,1.75);
\node at (5.35,1.2) {\scriptsize 入};
\node at (7.28,1.2) {\scriptsize 反};
\draw[->,violet,thick] (6.55,0.28) -- (8.05,0.28);
\node at (8.28,0.28) {\scriptsize $\vec M$};

% --- (c) transverse ---
\node at (11.9,2.25) {\textbf{(c) 横向 transverse}};
\draw[thick,violet] (10.0,0.55) -- (13.9,0.55);
\draw[->,red!65!black] (10.1,1.75) -- (11.3,0.56);
\draw[->,red!65!black] (11.3,0.56) -- (12.5,1.75);
\node at (10.35,1.2) {\scriptsize 入};
\node at (12.28,1.2) {\scriptsize 反};
\draw[violet,thick] (11.95,0.28) circle (0.12);
\draw[violet,thick] (11.885,0.215) -- (12.015,0.345);
\draw[violet,thick] (11.885,0.345) -- (12.015,0.215);
\node at (12.28,0.28) {\scriptsize $\vec M$};
\end{tikzpicture}
\end{center}

{\small\color{inkgray}图 6：MOKE 三种几何。紫色单箭头与 $\otimes$ 是磁化方向（$\otimes$ = 垂直纸面向里）。
(a) $\vec M$ 垂直膜面；(b) $\vec M$ 在面内、平行入射面；(c) $\vec M$ 在面内、垂直入射面。
判据是 $\vec M$ 相对「膜面法线」与「入射面」这两个方向的位置，与光从哪一侧来无关。}

{\small\color{inkgray}图 6：MOKE 三种几何。$\otimes$ 表示磁化垂直纸面向里（在面内、垂直于入射面）。
判据看 $\vec M$ 相对「表面法线」与「入射面」两个方向的位置。}

\begin{center}
\small
\tabtop
\begin{tabular}{@{}p{3.2cm}p{5.0cm}p{5.8cm}@{}}
\thead{几何} & \thead{磁化方向} & \thead{测什么} \\
\tabmid
极向 polar & 垂直样品表面 & 面外磁化；正入射也可测 \\
\tabmid
纵向 longitudinal & 面内、平行入射面 & 最常见的薄膜磁滞回线 \\
\tabmid
横向 transverse & 面内、垂直入射面 & 只测强度变化，不测旋转 \\
\end{tabular}
\tabbot
\normalsize
\end{center}

\subhdr{13.3 复克尔转角（课件 p62--p63）}

MOKE 测量的是「$p$ 偏振入射时出现的 $s$ 分量」——反射系数矩阵有了非对角元：
\begin{equation}
\label{eq:moke}
\Phi=\theta+i\psi=\frac{r_{ps}}{r_{pp}} ,
\qquad
\Phi\propto \varepsilon_{xy}\ (\text{即}\propto Q),
\end{equation}
实部是旋转、虚部是椭圆率。$r_{ps}\neq0$ 的根源就是 $\varepsilon_{xy}$：没有磁化时，$p$ 与 $s$ 互不混合。

极向几何下 $r_{ps}$ 的具体形式（课件 p63）是
\begin{equation}
\theta_p=\mathrm{Im}\!\left[
\frac{\eta^2}{\eta^2-\sin^2\theta_i}\cdot
\frac{\mp\sin\theta_i\tan\theta_i\,(\eta^2-1)}{\eta^2-\tan^2\theta_i}\cdot Q
\right],
\qquad
\psi_p=\mathrm{Re}\,[\,\cdots\,],
\end{equation}
其中 $\eta=\eta_2/\eta_1$ 是磁光层与环境的复折射率之比，$\theta_i$ 是入射角，$Q=g/\varepsilon$。
\textbf{这个式子的用途是判断趋势，不是拿来背}：$Q\to0$ 时 $\theta_p\to0$（无磁化无信号）；
$\eta$ 与 $\theta_i$ 决定响应的大小与符号。考试若考，问到的是「$\Phi$ 的实部虚部各是什么」与「$Q$ 与什么成正比」。

\subhdr{13.4 纵向与横向（课件 p64--p65）}

\textbf{纵向 MOKE}：磁化在面内且平行入射面。$\theta_L$ 有与 $\theta_p$ 同构的表达式，但多一个
$\sin\theta_i\eta^2/\sqrt{\eta^2-\sin^2\theta_i}$ 的因子。\textbf{关键性质：正入射时 $\theta_L=0$}——
因为正入射没有"入射面"，纵向与极向无法区分，而只有垂直膜面的磁化才能在正入射下起作用。

\textbf{横向 MOKE}：磁化在面内且垂直入射面。此时 $p$ 分量受影响、$s$ 分量不受影响，
所以\textbf{偏振态不变、只改反射强度}：
\begin{equation}
\delta T=\frac{I_M-I_{-M}}{I_0}
=-\mathrm{Im}\!\left[\frac{4\tan\theta_i\,\eta^2}{\eta^2-1}\cdot
\frac{\eta^2}{\eta^2-\tan^2\theta_i}\cdot Q\right].
\end{equation}
$I_M$ 与 $I_{-M}$ 是相反磁化方向下的反射强度。$p$ 偏振斜入射时这个效应最大。

\subhdr{13.5 Faraday 与 MOKE 的分工（课件 p77）}

\begin{center}
\small
\tabtop
\begin{tabular}{@{}p{3.0cm}p{5.4cm}p{5.2cm}@{}}
\thead{项目} & \thead{Faraday} & \thead{MOKE} \\
\tabmid
几何 & 透射 & 反射 \\
\tabmid
测量量 & 透射偏振面的旋转 & 反射的旋转 + 椭圆率 \\
\tabmid
探测深度 & 整个样品厚度 & 表面（金属约 $20\ \mathrm{nm}$） \\
\tabmid
典型应用 & 隔离器、电流传感器 & 磁畴成像、薄膜表征 \\
\end{tabular}
\tabbot
\normalsize
\end{center}

一句话：\textbf{要看体内用 Faraday，要看表面用 MOKE}。

%=====================================================================
\hdr{§14}{Voigt/Cotton--Mouton 与圆二色性}

\subhdr{14.1 磁线双折射：Voigt 与 Cotton--Mouton（课件 p66--p67）}

前面所有磁光效应都是\emph{圆}双折射或圆二色性（需要 $\vec M\parallel\vec k$）。若把磁化方向转到
\emph{垂直}于传播方向，得到的是\emph{线}双折射，而且它是\emph{二次}的：

\begin{center}
\small
\tabtop
\begin{tabular}{@{}p{3.6cm}p{4.4cm}p{5.6cm}@{}}
\thead{名称} & \thead{介质} & \thead{几何} \\
\tabmid
Voigt 效应（1902） & 气体、蒸气 & $\vec B\perp\vec k$，线双折射 \\
\tabmid
Cotton--Mouton 效应（1907） & 液体 & 同上，同一现象的液体版 \\
\end{tabular}
\tabbot
\normalsize
\end{center}

\textbf{怎么理解它的二次性}：磁化沿横向时，把磁化反号等价于绕传播轴转 $180^\circ$，
而线双折射（$\vec E\parallel\vec M$ 与 $\vec E\perp\vec M$ 的相位差）在这个转动下不变——
所以响应只能是 $M$ 的\emph{偶}函数，最低阶是 $M^2$。这与 §2 的对称性论证是同一套逻辑：
\textbf{纵向效应（Faraday）是 $M$ 的奇函数，横向效应是偶函数}。

定量上（课件 p67）：
\begin{equation}
\Delta n_{\text{LM}}\propto I_\parallel-I_\perp,\qquad \text{且}\ \Delta n_{\text{LM}}\sim M^2 ,
\end{equation}
$I_\parallel$、$I_\perp$ 是磁化平行/垂直于偏振时的反射或透射强度。介质表现得像一块光轴沿 $\vec M$ 的单轴晶体。
量级上它比 Faraday 弱得多（$M^2$ 且系数小），主要用于气体与液晶的基础测量。

\subhdr{14.2 圆二色性 CD（课件 p70--p72）}

CD 是 RCP 与 LCP 的\emph{吸收}差——圆双折射的"吸收版"：
\begin{equation}
\label{eq:cd}
CD=\Delta A=A_{LCP}-A_{RCP}=(\alpha_- -\alpha_+)\,d ,
\end{equation}
$\alpha_\pm$ 是两个圆偏振的吸收系数，$d$ 是光程。$\Delta A\neq0$ 时，一个圆分量被削得更多，
出射光变成\emph{椭圆}偏振（这正是"二色性"的字面意思）。

\begin{center}
\small
\tabtop
\begin{tabular}{@{}p{3.6cm}p{2.6cm}p{7.4cm}@{}}
\thead{量} & \thead{符号} & \thead{含义} \\
\tabmid
吸光度差 & $\Delta A$ & 原始 CD 信号 \\
\tabmid
摩尔椭圆率 & $[\theta]$ & 对浓度与光程归一（$\mathrm{deg\cdot cm^2\cdot dmol^{-1}}$） \\
\tabmid
不对称因子 & $g=\Delta A/A$ & 手性强弱，典型 $10^{-5}$--$10^{-3}$ \\
\tabmid
椭圆率角 & $\theta$（mdeg） & $\tan\theta=$ 椭圆短轴/长轴 \\
\end{tabular}
\tabbot
\normalsize
\end{center}

\textbf{物理来源}：介质必须对两个圆偏振"区别对待"，两种途径——
（i）\textbf{手性}（缺镜面对称），如蛋白质、DNA、糖；（ii）\textbf{磁化}（MCD），连非手性样品也有。
微观机制分别是电偶极–磁偶极干涉（自然 CD）与 Zeeman 分裂（MCD），结果都是复折射率不等：
\begin{equation}
\tilde n_+\neq\tilde n_-\quad\Longrightarrow\quad \Delta n\neq0\ (\text{CB})\ \text{且}\ \Delta\alpha\neq0\ (\text{CD}).
\end{equation}

\textbf{仪器（课件 p72）}：光源 → 滤光轮（选波长）→ 偏振器 → 光弹调制器或 $\lambda/4$ 片
（交替产生 LCP 与 RCP）→ 样品池 → 探测器。测的是两个圆偏振下的吸收差，信噪比靠调制与锁相放大保证。

\subhdr{14.3 CB 与 CD 的对照（课件 p73）}

\begin{center}
\small
\tabtop
\begin{tabular}{@{}p{3.0cm}p{5.4cm}p{5.2cm}@{}}
\thead{方面} & \thead{圆双折射 CB} & \thead{圆二色性 CD} \\
\tabmid
差别在 & RCP 与 LCP 的\emph{相速}（折射率） & RCP 与 LCP 的\emph{吸收}（衰减） \\
\tabmid
物理量 & $\Delta n=n_+-n_-$（实部） & $\Delta\alpha=\alpha_+-\alpha_-$（虚部） \\
\tabmid
对光的作用 & 偏振面\emph{旋转} & 偏振态变\emph{椭圆} \\
\tabmid
联系 & \multicolumn{2}{l}{Kramers--Kronig 关系：一个频率上有 CD，附近必有 CB} \\
\end{tabular}
\tabbot
\normalsize
\end{center}

这张表把整讲的磁光部分收口了：\textbf{同一个复数的 $\varepsilon_{xy}$，实部给 CB（旋转）、虚部给 CD（椭圆）}；
Faraday（§11）是它的透射版，MOKE（§13）是它的反射版，Voigt（§14.1）是横向的线双折射版。

%=====================================================================
\hdr{§15}{材料、增强手段与应用}

\subhdr{15.1 磁光材料的四类（课件 p69）}

\begin{center}
\small
\tabtop
\begin{tabular}{@{}>{\raggedright}p{3.2cm}>{\raggedright}p{5.6cm}>{\raggedright}p{3.6cm}>{\raggedright\arraybackslash}p{3.2cm}@{}}
\thead{类型} & \thead{代表材料} & \thead{工艺} & \thead{适用条件} \\
\tabmid
铁磁金属 & Fe、Co、Ni、Fe$_{20}$Ni$_{80}$、Ag/Co、Tb$_{18}$Co$_{82}$ & 磁控溅射 & 室温；可见--近红外 \\
\tabmid
磁性电介质 & YIG、Bi:YIG、Ce:YIG、BIG、Fe$_3$O$_4$、CoFe$_2$O$_4$、EuS & 溅射、LPE、PLD & 室温；可见--近红外 \\
\tabmid
磁性半导体 & 石墨烯、InSb、InAs、Cd$_x$Mn$_{1-x}$Te & 外延 & 室温；太赫兹 \\
\tabmid
稀土化合物与拓扑材料 & EuS、EuO、CrI$_3$、MnBi、EuCd$_2$As$_2$ & 物理气相沉积 & 低温；可见--近红外 \\
\end{tabular}
\tabbot
\normalsize
\end{center}

选材料的顺序与 §6.3 的电光材料类似：\textbf{一、工作波段} → \textbf{二、Verdet 常数 / 磁光常数}
→ \textbf{三、能否在室温工作}（第三、四类多半要低温）→ \textbf{四、工艺兼容性}（能否集成到芯片上）。
铁磁金属的磁光信号大但吸收也大（只适合反射式的 MOKE 或薄膜）；器件（隔离器）主要靠磁性电介质。

\subhdr{15.2 为什么天然材料的磁光效应弱，怎么增强（课件 p78）}

\textbf{弱的根源}：磁光效应来自自旋–轨道耦合，而自旋–轨道耦合相对交换劈裂与光跃迁能量是小量
（它是相对论修正）。结果两个圆偏振的折射率差通常只有 $10^{-5}$--$10^{-3}$，单位长度的转角很小。

\textbf{五条增强途径}：
\begin{enumerate}
\item \textbf{强自旋–轨道耦合的离子}：Ce$^{3+}$、Bi$^{3+}$、Tb$^{3+}$ 掺进石榴石，转角显著增大；
\item \textbf{交换增强}：YIG 里交换场作用在稀土离子上，响应远超外场单独能给出的；
\item \textbf{共振增强}：工作在跃迁附近（如 Ce:YIG 的 Ce$^{3+}\to$Fe$^{3+}$ 电荷转移），转角可升几个数量级；
\item \textbf{多层膜与纳米结构}：等离子体或介质共振在局部增强相互作用；
\item \textbf{强磁场 / 饱和磁化}：效应正比于磁化，先饱和再谈别的。
\end{enumerate}
这就是实用器件用 Ce:YIG、TGG 这类"工程化"材料、而不是简单铁磁金属的原因。

\subhdr{15.3 磁光存储（课件 p75）}

磁光盘（MO 盘）把 §13 的读法与热效应串起来：\textbf{写}——激光脉冲把一点加热到居里温度以上
（矫顽力降到零），外场按数据位设定磁化方向，冷却后固定；\textbf{读}——极向 MOKE，
读出光的偏振面按该点磁化方向转 $+\theta$ 或 $-\theta$，差分探测读出"0/1"（转角只有零点几度，
所以对信噪比要求高）。历史：1957 年首次磁光记录（MnBi 薄膜），1980--1990 年代商品化。

\subhdr{15.4 器件材料与传感}

隔离器/环形器要的是「每特斯拉每米转多少度」：由 $L=\theta/(\vdt B)$ 可见 $\vdt$ 大 $4$ 倍、长度就短 $4$ 倍，
所以选材就是「在目标波段找 $\vdt$ 最大、吸收最小的」——TGG 用于可见--近红外，Ce:YIG 用于集成（薄膜）器件。

\textbf{光纤电流传感器}：绕被测导线一圈的光纤环里，Faraday 转角正比于环内总电流（安培环路定理）。
技巧是\textbf{让同一束光走两趟做共模抑制}——Faraday 往返是 $2\theta$，而应力双折射这类互易效应往返抵消，
信号加倍、干扰抵消（又一次用到 §12）。

\textbf{一句话收口}：电光效应给「改相位」的旋钮，磁光效应给「改方位、且改不回来」的旋钮；
两个旋钮加偏振器，就是调制器、Q 开关、隔离器三件主力器件。

%=====================================================================
\needspace{16\baselineskip}
\hdr{§16}{作业（课件 p82）}

{\small\color{inkgray}三题题干如下（不含关键步骤与答案）；中文详解见附录 C，英文可抄写版见附录 D。}

\practice{P1}{课件 p82 Problem 1，第 1 题：Pockels 盒的相位延迟}{一块 KDP 晶体的普克尔斯盒，电光系数 $r=10.6\times10^{-12}\ \mathrm{m/V}$，
晶体长 $3\ \mathrm{cm}$，两端加 $5\ \mathrm{kV}$ 电压（孔径即电极间距 $d=5\ \mathrm{mm}$），
取 $n=1.5$、$\lambda=633\ \mathrm{nm}$。(a) 求感生的相位延迟 $\Delta\phi$；(b) 求这块晶体的半波电压 $V_\pi$。}
{\textbf{几何判定}：给出 $L$ 与 $d$，电压加在 $d$ 上 → 横向几何，$E=V/d$；
相位延迟用 $\Gamma=\frac{2\pi}{\lambda}n^3rE\,L$（注意两个模之差是 $n^3rE$，不带 $\frac12$）；
$V_\pi$ 与 $\Gamma$ 的自洽关系 $\Gamma=\pi V/V_\pi$}{ans:P1}

\practice{P2}{课件 p82 Problem 2，第 2 题：Faraday 旋转与往返}{波长 $\lambda=589\ \mathrm{nm}$ 的光束通过一根长 $L=10\ \mathrm{cm}$ 的玻璃棒，
棒置于产生 $B=0.5\ \mathrm{T}$ 磁场的螺线管中。玻璃的 Verdet 常数为 $\vdt=20\ \mathrm{rad/(T\cdot m)}$。
(a) 求偏振面的转角；(b) 若光两次穿过该棒（往返），总转角是多少？解释为什么结果不是"简单地翻倍"。}
{\textbf{公式}：$\theta=\vdt BL$（$L$ 要化成米）；
\textbf{非互易性}：旋转方向由 $\vec B$ 定，回程继续同向累加 → $2\theta$；
\textbf{措辞}：题目要解释的是「为什么不能按反向传播吃掉去程的直觉去算」，而不是「答案不是 $2\theta$」}{ans:P2}

\practice{P3}{课件 p82 Problem 3，第 3 题：Kerr 与 Pockels 的对照}{两块相同的电光晶体参与实验：晶体 A 中心对称，晶体 B 非中心对称。
两者都加 $E=10^{6}\ \mathrm{V/m}$ 的电场。
(a) 哪块显 Pockels 效应、哪块显 Kerr 效应？
(b) Kerr 效应的感生双折射为 $\Delta n=K\lambda E^2$，其中 $K=2\times10^{-14}\ \mathrm{m/V^2}$、$\lambda=633\ \mathrm{nm}$，求 $\Delta n$。
(c) 若晶体长 $L=1\ \mathrm{cm}$，求两个偏振分量之间的相位差 $\Delta\phi$。}
{\textbf{对称性判据}：非中心对称才有线性项；
\textbf{量纲}：$K\lambda E^2$ 三个单位的乘积要无量纲；
\textbf{换算}：$\Delta\phi=\frac{2\pi}{\lambda}\Delta n\,L$，注意此处 $\lambda$ 会与 $\Delta n$ 里的 $\lambda$ 相消}{ans:P3}

\clearpage
\appendix

%=====================================================================
\hdr{附录 A}{公式速查}

\begin{center}
\small
\tabtop
\begin{tabular}{@{}p{6.2cm}p{8.8cm}@{}}
\thead{公式} & \thead{条件与说明} \\
\tabmid
$\Gamma=\dfrac{2\pi}{\lambda}\Delta n\,L$ & 任意双折射的相位延迟；$\lambda$ 为真空波长 \\
\tabmid
$\Delta n=K\lambda E^2$ & Kerr：所有材料；$K$ 单位 $\mathrm{m/V^2}$ \\
\tabmid
$\Gamma=2\pi KLE^2=2\pi KL\left(\dfrac{V}{d}\right)^2$ & Kerr 盒；$\lambda$ 已消去 \\
\tabmid
$\Delta n=-\dfrac12n^3rE$ & Pockels：\emph{单个}模的折射率变化（非中心对称材料） \\
\tabmid
$\Delta n_{\text{两模}}=n^3rE$ & Pockels：两个正交模之差；$\frac12$ 在此抵消 \\
\tabmid
$\Gamma_\parallel=\dfrac{2\pi}{\lambda}n^3rV$ & 纵向：与 $L$、$d$ 都无关 \\
\tabmid
$\Gamma_\perp=\dfrac{2\pi}{\lambda}n^3rV\dfrac{L}{d}$ & 横向：靠 $L/d$ 降电压 \\
\tabmid
$V_\pi=\dfrac{\lambda}{2n^3r}\cdot\dfrac{d}{L}$ & 纵向取 $d/L=1$ \\
\tabmid
$\Gamma=\pi\dfrac{V}{V_\pi}$ & 任何 Pockels 盒的通用换算 \\
\tabmid
$T=\sin^2\dfrac{\Gamma}{2}=\sin^2\dfrac{\pi V}{2V_\pi}$ & 正交偏振器之间 \\
\tabmid
$T\approx\dfrac12+\dfrac{\pi}{2}\dfrac{v}{V_\pi}$ & Q 点（$V=V_\pi/2$）附近的小信号线性化 \\
\tabmid
$\theta=[\alpha]\,c\,L$ & 自然旋光（Biot 定律）；糖度计 \\
\tabmid
$\theta=\dfrac{\pi}{\lambda}(n_--n_+)L$ & 圆双折射引起的转角（互易） \\
\tabmid
$\theta=\vdt BL=\vdt\displaystyle\int\vec B\cdot\mathrm d\vec l$ & Faraday 旋转（非互易） \\
\tabmid
$\theta_{\text{往返}}=2\vdt BL$ & Faraday 双程；隔离器的原理 \\
\tabmid
$\Phi=\theta+i\psi=\dfrac{r_{ps}}{r_{pp}}$ & MOKE 的复克尔转角 \\
\tabmid
$CD=\Delta A=(\alpha_--\alpha_+)d$ & 圆二色性；$g=\Delta A/A$ \\
\tabmid
$\delta T=\dfrac{I_M-I_{-M}}{I_0}$ & 横向 MOKE 的相对强度变化 \\
\end{tabular}
\tabbot
\normalsize
\end{center}

%=====================================================================
\hdr{附录 B}{缺口清单与口径核对}

{\small\color{inkgray}本节记录两件事：课件跳过、而理解主线需要的环节（B.1）；课件原页里需要说明的口径问题（B.2）。
后者一律经复算或按定义核对，不是凭印象判定。}

\subhdr{B.1 课件跳过、本讲义补上的环节}

\begin{center}
\small
\tabtop
\begin{tabular}{@{}p{2.6cm}p{4.2cm}p{8.2cm}@{}}
\thead{课件位置} & \thead{跳过什么} & \thead{本讲义补在哪里、为什么必须补} \\
\tabmid
p22 & $\Delta n=-\frac12n^3rE$ 直接给出 & §5.1 从 $1/n'^2=1/n^2+rE$ 一阶展开推出来；并且说明 $\frac12$ 在算两模之差时抵消（§5.1 的框） \\
\tabmid
p26、p30 & 纵横两式直接给出 & §6.1 用 $E=V/L$ 与 $E=V/d$ 各推一次，看清 $L$ 为什么在纵向被消掉 \\
\tabmid
p30 & $T=\sin^2(\Gamma/2)$ 直接给出 & §7.2 用琼斯矢量四步推出（关键在检偏器投影的两个符号相反） \\
\tabmid
p2（计划第 3 项） & 自然旋光与糖度计在正课页缺失 & §9 补齐 Fresnel 圆双折射的推导；§12 的非互易性离不开这个对照 \\
\tabmid
p36 & 只列了 Q 开关，没有原理页 & §7.7 补到能读懂器件：为什么用 $V_\pi/2$、为什么往返累积 \\
\tabmid
p57、p59 & 只给张量与非对角元 & §11.2 说明「旋转看 Im、椭圆看 Re」，并把公式还原成可测量的 $\theta=\vdt BL$ \\
\tabmid
p61--p65 & 三种 MOKE 只给结论与 $\theta_p$、$\theta_L$ 的公式 & §13 给出统一框架 $\Phi=r_{ps}/r_{pp}$，并说明正入射下 LMOKE 为零的原因 \\
\end{tabular}
\tabbot
\normalsize
\end{center}

\subhdr{B.2 口径核对（六条）}

\textbf{1.（p59）Faraday 公式缺少长度与折射率因子。}课件写作
$\theta=-\frac{\omega}{2c}\mathrm{Im}(\varepsilon_{xy}/\varepsilon)$、$\psi=\frac{\omega}{2c}\mathrm{Re}(\varepsilon_{xy}/\varepsilon)$。
按定义 $n_\pm^2=\varepsilon\mp i\varepsilon_{xy}$ 解出 $n_\pm\approx n\left(1\mp i\varepsilon_{xy}/2\varepsilon\right)$，
$n_+-n_-=-i\varepsilon_{xy}/n$；再由 $\theta=\frac{\pi}{\lambda}(n_--n_+)L$ 得
$\theta=-\frac{\omega L}{2c}\cdot n\,\mathrm{Im}(\varepsilon_{xy}/\varepsilon)$（同法得 $\psi$）。
所以课件两式是\emph{单位长度}的量，并且省掉了因子 $n$（若把 $\varepsilon$ 当绝对值则可吸收进系数）。
\textbf{结论}：结构与符号对，量纲要补 $L$；作业一律用 $\theta=\vdt BL$。

\textbf{2.（p53/p56 与 p59）术语缩写。}p54 正确地说 MO 效应表现为 MCB（磁圆双折射）与 MCD（磁圆二色性）；
p59 结尾却把 Faraday 效应称作 "Magnetic Circular Birefringence (MCD)"——缩写与名称对不上：
圆\emph{双折射}是 MCB，MCD 是圆\emph{二色性}（吸收差）。按定义应为 MCB。

\textbf{3.（p28）KDP 的 $V_\pi=7.5\ \mathrm{kV}$ 对应哪个波长。}由 $V_\pi=\lambda/(2n^3r)$，取 $n=1.51$、$r=10.5\ \mathrm{pm/V}$：
$633\ \mathrm{nm}$ 得 $8.75\ \mathrm{kV}$，$546\ \mathrm{nm}$ 得 $7.55\ \mathrm{kV}$。
表里的 $7.5\ \mathrm{kV}$ 与 $546\ \mathrm{nm}$ 吻合，标注的 $633\ \mathrm{nm}$ 应是沿用习惯值。表是"典型器件值"，
与 p43--p45 例题的 $n$、$r$ 不是同一组，不能混用。

\textbf{4.（p60）MOKE 的"沿传播方向的磁化分量"。}该说法只在极向几何成立。纵向 MOKE 正比于面内、平行入射面的分量，
横向 MOKE 正比于面内、垂直入射面的分量；正入射时纵向与横向都不敏感（$\theta_L=0$）。

\textbf{5.（p57）"Absorption → Voigt Effect"。}Voigt/Cotton--Mouton 是磁\emph{线双折射}（相位差），
不是吸收效应；与吸收对应的是 MCD。该页的第三种分类应改为"横向几何 → Voigt 效应"。

\textbf{6.（p13 与 p46）两个 $K$ 差 100 倍。}p13 表里硝基苯 $K\sim2.2\times10^{-12}\ \mathrm{m/V^2}$，
p46 例题用 $K=2.2\times10^{-14}$。两者不是同一材料（后者接近 CS$_2$ 的量级），做题时按题给值算，
不要拿 p13 的表去替换题目里的 $K$。

%=====================================================================
\hdr{附录 C}{作业三题详解（中文）}

\answerstart{P1}{课件 p82 Problem 1}{KDP，$r=10.6\times10^{-12}\ \mathrm{m/V}$，$L=3\ \mathrm{cm}$，
$V=5\ \mathrm{kV}$，$d=5\ \mathrm{mm}$，$n=1.5$，$\lambda=633\ \mathrm{nm}$。
(a) 求 $\Delta\phi$；(b) 求 $V_\pi$。}

\textbf{第一步：判定几何，写出 $E$。}题目给的 $d=5\ \mathrm{mm}$ 是孔径（电极间距），电压加在它上面，
所以是\textbf{横向}配置：
\begin{equation}
E=\frac{V}{d}=\frac{5\times10^{3}}{5\times10^{-3}}=1.0\times10^{6}\ \mathrm{V/m}.
\end{equation}
（若读成纵向，则 $L$ 根本不出现在结果里；而题目特意给了 $L$，也特意给了 $d$——这是判定为横向的线索。）

\textbf{第二步：$\Delta\phi$。}用 $\Gamma=\frac{2\pi}{\lambda}n^3rE\,L$，注意两模之差是 $n^3rE$（§5.1）：
\begin{align}
\Gamma&=\frac{2\pi}{633\times10^{-9}}\times(1.5)^3\times10.6\times10^{-12}\times1.0\times10^{6}\times0.03 \nonumber\\
&=9.926\times10^{6}\times3.375\times10.6\times10^{-12}\times3.0\times10^{4}\nonumber\\
&\approx10.65\ \mathrm{rad}.
\end{align}
换成角度与 $\pi$ 的倍数：$\Gamma\approx 10.65\times\frac{180^\circ}{\pi}\approx610^\circ$，
或 $\Gamma\approx3.39\pi$。

\textbf{第三步：$V_\pi$。}两条等价路线，用哪条都行。

路线一（直接代公式）：
\begin{equation}
V_\pi=\frac{\lambda}{2n^3r}\cdot\frac{d}{L}
=\frac{633\times10^{-9}}{2\times3.375\times10.6\times10^{-12}}\times\frac{5\times10^{-3}}{3\times10^{-2}}
\approx1.47\times10^{3}\ \mathrm{V}.
\end{equation}
路线二（用自洽关系 $\Gamma=\pi V/V_\pi$）：
\begin{equation}
V_\pi=\frac{\pi V}{\Gamma}=\frac{\pi\times5000}{10.65}\approx1.47\times10^{3}\ \mathrm{V}.
\end{equation}
两条路线必须一致，这是\textbf{自检}：$10.65\approx\pi\times5000/1474$\ \checkmark。

\textbf{答}：(a) $\Delta\phi\approx1.07\times10^{1}\ \mathrm{rad}\ (\approx610^\circ\approx3.4\pi)$；
(b) $V_\pi\approx1.5\ \mathrm{kV}$。

\textbf{点评}。(i) 施加的 $5\ \mathrm{kV}$ 已经是 $V_\pi$ 的 $3.4$ 倍，所以延迟超过了三个半波——
这说明这块晶体"太长了"：同样的电压若只想要半波延迟，$L$ 应缩短到约 $0.9\ \mathrm{cm}$。
(ii) 若把题目读成纵向配置，则 $\Delta\phi=1.78\ \mathrm{rad}$、$V_\pi=8.85\ \mathrm{kV}$；
两种读法的差别正是「$E=V/L$ 还是 $E=V/d$」，写作业时把选择写出来即可。
(iii) 实际器件工作在 $V_\pi/2\approx0.74\ \mathrm{kV}$（Q 点，§7.3），而不是 $V_\pi$。

%---------------------------------------------------------------------
\answerstart{P2}{课件 p82 Problem 2}{玻璃棒 $L=10\ \mathrm{cm}$、$B=0.5\ \mathrm{T}$、$\vdt=20\ \mathrm{rad/(T\cdot m)}$，
$\lambda=589\ \mathrm{nm}$。(a) 转角；(b) 往返总转角并解释。}

\textbf{(a)} 直接代 $\theta=\vdt BL$（$L$ 先化成米）：
\begin{equation}
\theta=20\times0.5\times0.10=1.0\ \mathrm{rad}\;\approx\;57.3^\circ .
\end{equation}
（$\mathrm{rad/(T\cdot m)}\times\mathrm{T}\times\mathrm{m}=\mathrm{rad}$，量纲自洽。）
题目给的 $\lambda=589\ \mathrm{nm}$ 在公式里不出现——$\lambda$ 的影响藏在 $\vdt$ 里（$\vdt$ 随波长效增），
所以答题时说明「$\vdt$ 已按该波长取值」即可。

\textbf{(b)} 往返总转角：
\begin{equation}
\theta_{\text{往返}}=2\vdt BL=2.0\ \mathrm{rad}\approx114.6^\circ .
\end{equation}

\textbf{为什么不是「简单地翻倍」这句要这样理解。}物理上它\emph{确实}是两倍——但为什么是两倍而不是抵消，
需要说清机制：

\begin{enumerate}
\item 旋转的\emph{方向}由 $\vec B$ 决定，不由光的传播方向决定。反射只把 $\vec k$ 反号，$\vec B$ 不变。
\item 介质的介电张量由 $\vec B$ 固定，张量不变 → 两个圆偏振模的折射率 $n_\pm$ 在回程与去程\textbf{相同}。
\item 用 $\theta=\vdt\int\vec B\cdot\mathrm d\vec l$ 看：去程 $\mathrm d\vec l\parallel\vec B$，回程 $\mathrm d\vec l\parallel-\vec B$——
      但旋转的符号也随之定义了"相对 $\vec B$ 的正方向"，两趟在实验室坐标里是同向累加。
      严格的相位计算（§12.2 三步）给出相对相位差 $2(k_+-k_-)L$，即 $2\theta$。
\end{enumerate}

\textbf{对照}：若换成糖溶液或石英（自然旋光），回程的 $n_\pm$ 交换，转角反号，往返合计为 $0$。
所以「不能简单地按反向传播就抵消」才是这句话的要点；而「2 倍」是非互易性的\emph{结果}，不是巧合。

若想让答案无懈可击，加一句英文：
\emph{The doubling is not a coincidence: the rotation sense is locked to $\mathbf{B}$, so the return pass
adds $+\theta$ again in the laboratory frame (total $2\theta$), whereas a reciprocal rotator would give $0$.}

\textbf{答}：(a) $\theta=1.0\ \mathrm{rad}=57.3^\circ$；(b) $2\theta=2.0\ \mathrm{rad}=114.6^\circ$，因为 Faraday 效应非互易：
决定旋向的是 $\vec B$ 而不是 $\vec k$，回程继续同向累加。

%---------------------------------------------------------------------
\answerstart{P3}{课件 p82 Problem 3}{A 中心对称、B 非中心对称，$E=10^{6}\ \mathrm{V/m}$。
(a) 谁显 Pockels、谁显 Kerr；(b) $\Delta n=K\lambda E^2$，$K=2\times10^{-14}$、$\lambda=633\ \mathrm{nm}$，求 $\Delta n$；
(c) $L=1\ \mathrm{cm}$，求 $\Delta\phi$。}

\textbf{(a)} 由 §2 的判据：中心对称介质的 $\Delta n(E)$ 只能是 $E$ 的偶函数，奇次项被禁。
\begin{itemize}
\item \textbf{晶体 A（中心对称）}：Pockels 被禁，\textbf{只显 Kerr 效应}；
\item \textbf{晶体 B（非中心对称）}：\textbf{显 Pockels 效应}（线性项允许，通常远大于二次项）；Kerr 效应也同时存在，
      只是被线性项掩盖。
\end{itemize}

\textbf{(b)}
\begin{equation}
\Delta n=K\lambda E^2=2\times10^{-14}\times633\times10^{-9}\times\left(10^{6}\right)^2
=2\times10^{-14}\times6.33\times10^{5}\approx1.27\times10^{-8}.
\end{equation}
量纲自检：$\frac{\mathrm{m}}{\mathrm{V^2}}\cdot\mathrm{m}\cdot\frac{\mathrm{V^2}}{\mathrm{m^2}}=1$（无量纲），\checkmark。

\textbf{(c)}
\begin{equation}
\Delta\phi=\frac{2\pi}{\lambda}\Delta n\,L
=\frac{2\pi}{633\times10^{-9}}\times1.27\times10^{-8}\times0.01
\approx1.26\times10^{-3}\ \mathrm{rad}.
\end{equation}
换成角度约 $0.072^\circ$。

\textbf{一个更快的算法（建议写在卷面上作为自检）}：把 $\Delta n=K\lambda E^2$ 直接代进 $\Delta\phi$，$\lambda$ 相消：
\begin{equation}
\Delta\phi=2\pi KLE^2=2\pi\times2\times10^{-14}\times0.01\times10^{12}\approx1.26\times10^{-3}\ \mathrm{rad},
\end{equation}
两式必须给出同一个数\ \checkmark。这个消去还说明：\textbf{Kerr 盒的相位延迟与波长无关}（在 $K$ 缓变的近似下）。

\textbf{点评}。(i) $\Delta n\sim10^{-8}$ 比 Pockels 在同样场强下的 $\sim1.8\times10^{-5}$ 小约 $10^3$ 倍
（$\frac12n^3rE$ 取 $n=1.5$、$r=10.6\ \mathrm{pm/V}$），这是「能用 Pockels 就不用 Kerr」的定量依据。
(ii) 相位差只有 $10^{-3}\ \mathrm{rad}$，在正交检偏器下对应的强度变化 $\sim\sin^2(5\times10^{-4})\sim10^{-7}$——
几乎测不出来。要做出可用的 Kerr 器件，必须把 $E$ 提高一到两个数量级、把 $L$ 加长，或换更强的 Kerr 液体。

\textbf{答}：(a) A 显 Kerr、B 显 Pockels；(b) $\Delta n\approx1.3\times10^{-8}$；
(c) $\Delta\phi\approx1.3\times10^{-3}\ \mathrm{rad}$（约 $0.07^\circ$）。

%=====================================================================
\hdr{附录 D}{English solution sheet · home assignment (courseware p82)}

\callout{\textbf{How to use this appendix.} The three problems on courseware p82 are to be handed in
\emph{handwritten in English}. Copy the working below; keep the four parts the instructor asks for:
\textbf{Given}, \textbf{Formulas used}, \textbf{step-by-step calculation}, \textbf{final answers with units}.
Sign your name in English and put the sheet on the desk before the lecture starts. The Chinese commentary
stays in appendix C; this appendix contains only what goes on the paper.}

\subhdr{D.1 Question 1 — Pockels cell phase retardation}

\callout{\textbf{Problem 1.} A Pockels cell is made from a potassium dihydrogen phosphate (KDP) crystal with
an electro-optic coefficient of $r=10.6\times10^{-12}\ \mathrm{m/V}$. The crystal is $3\ \mathrm{cm}$ long, and a
voltage of $5\ \mathrm{kV}$ is applied across it (aperture $d=5\ \mathrm{mm}$). Assume $n=1.5$ and
$\lambda=633\ \mathrm{nm}$. (a) Calculate the induced phase retardation $\Delta\phi$. (b) Determine the
half-wave voltage $V_\pi$ for this crystal.}

\textbf{Given}\quad KDP, $r=10.6\times10^{-12}\ \mathrm{m/V}$, $n=1.5$, $\lambda=633\ \mathrm{nm}$,
crystal length $L=3.0\times10^{-2}\ \mathrm{m}$, electrode spacing $d=5.0\times10^{-3}\ \mathrm{m}$,
applied voltage $V=5.0\times10^{3}\ \mathrm{V}$. The field is applied across $d$ and the light travels along
$L$, so the configuration is \emph{transverse} and $E=V/d$.

\textbf{Formulas used}\quad
$E=\dfrac{V}{d}$;\quad $\Delta\phi=\dfrac{2\pi}{\lambda}\,n^{3}r\,E\,L$\quad
(the two eigen-polarizations differ by $n^{3}rE$, not by $\tfrac12 n^{3}rE$);\quad
$V_\pi=\dfrac{\lambda}{2n^{3}r}\cdot\dfrac{d}{L}$;\quad $\Delta\phi=\pi\dfrac{V}{V_\pi}$.

\textbf{(a)}
\begin{equation}
E=\frac{5.0\times10^{3}}{5.0\times10^{-3}}=1.0\times10^{6}\ \mathrm{V/m},
\end{equation}
\begin{equation}
\Delta\phi=\frac{2\pi}{633\times10^{-9}}\times1.5^{3}\times10.6\times10^{-12}\times1.0\times10^{6}\times3.0\times10^{-2}
\approx1.07\times10^{1}\ \mathrm{rad}.
\end{equation}
That is $\Delta\phi\approx610^\circ$, or about $3.4\pi$.

\textbf{(b)}
\begin{equation}
V_\pi=\frac{633\times10^{-9}}{2\times1.5^{3}\times10.6\times10^{-12}}\times\frac{5.0\times10^{-3}}{3.0\times10^{-2}}
\approx1.47\times10^{3}\ \mathrm{V}.
\end{equation}
Check with the other route: $V_\pi=\pi V/\Delta\phi=\pi\times5.0\times10^{3}/10.65\approx1.47\times10^{3}\ \mathrm{V}$.

\textbf{Answers}\quad (a) $\Delta\phi\approx1.1\times10^{1}\ \mathrm{rad}\ (\approx610^\circ\approx3.4\pi)$;\quad
(b) $V_\pi\approx1.5\times10^{3}\ \mathrm{V}\ (\approx1.5\ \mathrm{kV})$.

\subhdr{D.2 Question 2 — Faraday rotation and a double pass}

\callout{\textbf{Problem 2.} A beam of light with wavelength $\lambda=589\ \mathrm{nm}$ passes through a glass
rod of length $L=10\ \mathrm{cm}$ placed in a solenoid that produces a magnetic field $B=0.5\ \mathrm{T}$.
The Verdet constant of the glass is $\mathcal{V}=20\ \mathrm{rad/(T\cdot m)}$.
(a) Calculate the angle of rotation of the polarization plane.
(b) If the light passes through the rod twice (back and forth), what is the total rotation angle?
Explain why the result is not simply doubled.}

\textbf{Given}\quad $L=0.10\ \mathrm{m}$, $B=0.5\ \mathrm{T}$, $\mathcal{V}=20\ \mathrm{rad/(T\cdot m)}$
(taken at $\lambda=589\ \mathrm{nm}$), single pass in (a) and double pass in (b).

\textbf{Formulas used}\quad $\theta=\mathcal{V}BL$;\quad for a double pass
$\theta_{\text{round trip}}=2\mathcal{V}BL$ because the Faraday effect is \emph{non-reciprocal}.

\textbf{(a)}
\begin{equation}
\theta=\mathcal{V}BL=20\times0.5\times0.10=1.0\ \mathrm{rad}\approx57.3^\circ .
\end{equation}

\textbf{(b)}
\begin{equation}
\theta_{\text{round trip}}=2\,\mathcal{V}BL=2.0\ \mathrm{rad}\approx115^\circ .
\end{equation}

\textbf{Explanation.} The sense of the rotation is fixed by the magnetic field, not by the direction of
propagation. On reflection $\mathbf{k}$ reverses but $\mathbf{B}$ does not, so the dielectric tensor — and
therefore the two circular refractive indices $n_\pm$ — are unchanged; the relative phase $2\pi(n_+-n_-)L/\lambda$
accumulates a second time. In the laboratory frame the two passes add, giving $2\theta$, whereas a
\emph{reciprocal} rotator (natural optical activity in sugar solution or quartz) would give $-\theta$ on the
return pass and cancel to zero. The point of the question is that one must not apply the reciprocal
"it comes back, so it undoes itself" intuition.

\textbf{Answers}\quad (a) $\theta=1.0\ \mathrm{rad}=57.3^\circ$;\quad (b) $\theta_{\text{round trip}}=2.0\ \mathrm{rad}=115^\circ$,
because the rotation sense is set by $\mathbf{B}$ and not by $\mathbf{k}$, so the return pass adds another $+\theta$.

\subhdr{D.3 Question 3 — Kerr effect versus Pockels effect}

\callout{\textbf{Problem 3.} Two identical electro-optic crystals are used in an experiment. Crystal A is
centrosymmetric, and crystal B is non-centrosymmetric. An electric field $E=10^{6}\ \mathrm{V/m}$ is applied to
both. (a) Which crystal will exhibit the Pockels effect? Which will exhibit the Kerr effect?
(b) For the Kerr effect, the induced birefringence is $\Delta n=K\lambda E^{2}$, where
$K=2\times10^{-14}\ \mathrm{m/V^{2}}$ and $\lambda=633\ \mathrm{nm}$. Calculate $\Delta n$.
(c) If the crystal length is $L=1\ \mathrm{cm}$, calculate the phase difference $\Delta\phi$ between the two
polarization components.}

\textbf{Given}\quad crystal A centrosymmetric, crystal B non-centrosymmetric; $E=10^{6}\ \mathrm{V/m}$ for both;
in (b,c) $K=2\times10^{-14}\ \mathrm{m/V^{2}}$, $\lambda=633\ \mathrm{nm}$, $L=1.0\times10^{-2}\ \mathrm{m}$.

\textbf{Formulas used}\quad $\Delta n=K\lambda E^{2}$ (quadratic, allowed in every material);\quad
$\Delta\phi=\dfrac{2\pi}{\lambda}\Delta n\,L$;\quad symmetry rule: in a centrosymmetric medium
$\Delta n(+E)=\Delta n(-E)$, so all odd terms vanish.

\textbf{(a)} Crystal A (centrosymmetric) has inversion symmetry, so the linear term is forbidden:
it shows \textbf{only the Kerr effect}. Crystal B (non-centrosymmetric) allows the linear term: it shows the
\textbf{Pockels effect} (and the Kerr effect as well, but the linear term dominates).

\textbf{(b)}
\begin{equation}
\Delta n=K\lambda E^{2}=2\times10^{-14}\times633\times10^{-9}\times\left(10^{6}\right)^{2}
\approx1.27\times10^{-8}.
\end{equation}
Dimension check: $\dfrac{\mathrm{m}}{\mathrm{V^{2}}}\cdot\mathrm{m}\cdot\dfrac{\mathrm{V^{2}}}{\mathrm{m^{2}}}=1$, dimensionless.

\textbf{(c)}
\begin{equation}
\Delta\phi=\frac{2\pi}{\lambda}\Delta n\,L
=\frac{2\pi}{633\times10^{-9}}\times1.27\times10^{-8}\times1.0\times10^{-2}
\approx1.26\times10^{-3}\ \mathrm{rad}.
\end{equation}
Equivalently, substituting $\Delta n=K\lambda E^{2}$ first cancels $\lambda$:
$\Delta\phi=2\pi KLE^{2}\approx1.26\times10^{-3}\ \mathrm{rad}$ (about $0.072^\circ$).

\textbf{Answers}\quad (a) A: Kerr only; B: Pockels (linear) and Kerr;\quad
(b) $\Delta n\approx1.3\times10^{-8}$;\quad (c) $\Delta\phi\approx1.3\times10^{-3}\ \mathrm{rad}\ (\approx0.07^\circ)$.

\vspace{6pt}
{\small\color{inkgray}\textbf{Before you hand it in:} handwritten; Given, formulas, calculations, answers with
units; name signed in English; placed on the instructor's desk before the lecture starts. Each three-problem
set is worth 30 points.}

\end{document}

